\documentclass[UTF8, 12pt, fontset=fandol, oneside]{ctexart}
\usepackage{researchnotes}

\title{强化学习：同策略预测与函数近似}
\author{}
\date{}

\begin{document}

\maketitle

\section*{第二部分\quad 近似解法}

在本书第二部分中，我们把第一部分介绍的表格型方法加以扩展，使其能够应用于状态空间任意大的问题。
在许多我们希望应用强化学习的任务中，状态空间是组合式的、极其庞大的；例如，可能的摄像机图像数量远远大于宇宙中原子的数量。
在这种情形下，即使在时间和数据都无限的极限中，我们也不能指望找到一个最优策略或最优价值函数；相反，我们的目标是在有限计算资源下找到一个好的近似解。
在本书这一部分中，我们将考察这类近似解法。

大状态空间带来的问题，不只是大型表格所需的内存，还有准确填充这些表格所需的时间和数据。
在许多目标任务中，几乎每一个遇到的状态此前都从未出现过。
为了在这样的状态中作出合理决策，必须从过去遇到过的、在某种意义上与当前状态相似的不同状态进行泛化。
换句话说，关键问题就是\textbf{泛化}。
怎样才能把有限状态空间子集上的经验有用地泛化，从而在一个大得多的子集上产生良好近似？

幸运的是，从样例进行泛化已经得到了广泛研究，我们不需要为强化学习发明全新的方法。
在某种程度上，我们只需要把强化学习方法与已有的泛化方法结合起来。
我们所需的这类泛化通常称为\textbf{函数近似}，因为它从某个期望函数（例如价值函数）的样例出发，并试图由这些样例进行泛化，构造对整个函数的近似。
函数近似是\textbf{监督学习}的一个实例，而监督学习是机器学习、人工神经网络、模式识别和统计曲线拟合等领域研究的主要主题。
从理论上说，这些领域中研究过的任何方法，都可以在强化学习算法中充当函数近似器；不过在实践中，有些方法比另一些更容易适应这一角色。

带函数近似的强化学习会引入若干传统监督学习中通常不会出现的新问题，例如非平稳性、自举以及延迟目标。
我们将在这一部分的五章中依次介绍这些问题以及其他问题。
起初，我们把注意力限制在同策略训练上：第 9 章处理预测情形，即策略已给定，只近似它的价值函数；第 10 章处理控制情形，即寻找最优策略的一个近似。
带函数近似的离策略学习这一更具挑战性的问题将在第 11 章处理。
在这三章中，我们都必须回到第一性原理，重新考察学习目标，以便把函数近似纳入考虑。
第 12 章介绍并分析资格迹这一算法机制，它在许多情形中能显著改善多步强化学习方法的计算性质。
这一部分的最后一章考察一种不同的控制方法，即\textbf{策略梯度方法}；它们直接近似最优策略，而不必形成近似价值函数（不过如果同时近似价值函数和策略，它们可能会高效得多）。

\section*{第 9 章\quad 带近似的同策略预测}

在本章中，我们开始研究强化学习中的函数近似，方式是考虑它在由同策略数据估计状态价值函数中的应用；也就是说，用已知策略 $\pi$ 生成的经验来近似 $v_\pi$。
本章的新意在于，近似价值函数不再表示为一张表，而是表示为一种带权重向量 $\mathbf{w}\in\mathbb{R}^d$ 的参数化函数形式。
我们把给定权重向量 $\mathbf{w}$ 时状态 $s$ 的近似价值写作 $\hat{v}(s,\mathbf{w})\approx v_\pi(s)$。
例如，$\hat{v}$ 可以是状态特征的线性函数，其中 $\mathbf{w}$ 是特征权重向量。
更一般地，$\hat{v}$ 可以是由多层人工神经网络计算出的函数，其中 $\mathbf{w}$ 是所有层中连接权重组成的向量。
通过调整权重，网络可以实现范围很广的不同函数。
或者，$\hat{v}$ 也可以是由决策树计算出的函数，其中 $\mathbf{w}$ 是定义树的划分点和叶节点取值的所有数值。
通常，权重的个数（即 $\mathbf{w}$ 的维数）远小于状态数（$d\ll|\mathcal{S}|$），而改变一个权重会改变许多状态的估计价值。
因此，当某一个状态被更新时，这个改变会从该状态泛化出去，影响许多其他状态的价值。
这样的\textbf{泛化}使学习可能更强大，但也可能更难管理和理解。

也许有些出人意料，把强化学习扩展到函数近似，也会使它适用于部分可观测问题，即智能体无法获得完整状态的问题。
如果 $\hat{v}$ 的参数化函数形式不允许估计价值依赖于状态的某些方面，那么这些方面就相当于不可观测。
事实上，本书这一部分给出的所有关于使用函数近似方法的理论结果，同样适用于部分可观测的情形。
不过，函数近似做不到的是，用过去观测的记忆来增强状态表示。
第 17.3 节将简要讨论这类可能的进一步扩展。

\subsection*{9.1\quad 价值函数近似}

本书到目前为止讨论的所有预测方法，都被描述为对估计价值函数的更新；这些更新把特定状态处的价值向某个“备份值”或\textbf{更新目标}移动。
让我们把一次单独更新记为 $s\mapsto u$，其中 $s$ 是被更新的状态，$u$ 是 $s$ 的估计价值所要移向的更新目标。
例如，价值预测中的蒙特卡洛更新是 $S_t\mapsto G_t$，TD(0) 更新是 $S_t\mapsto R_{t+1}+\gamma\hat{v}(S_{t+1},\mathbf{w}_t)$，而 $n$ 步 TD 更新是 $S_t\mapsto G_{t:t+n}$。
在 DP（动态规划）策略评估更新中，
\[
s\mapsto \mathbb{E}_\pi\!\left[R_{t+1}+\gamma\hat{v}(S_{t+1},\mathbf{w}_t)\mid S_t=s\right],
\]
被更新的是任意状态 $s$；而在其他情形中，被更新的是实际经验中遇到的状态 $S_t$。

把每次更新解释为指定价值函数所期望的输入--输出行为的一个样例，这是很自然的。
在某种意义上，更新 $s\mapsto u$ 表示状态 $s$ 的估计价值应当更像更新目标 $u$。
到目前为止，实际更新一直很简单：状态 $s$ 的估计价值对应的表项只需向 $u$ 移动一小部分，所有其他状态的估计价值保持不变。
现在，我们允许任意复杂而精致的方法来实现更新，并且在 $s$ 处的更新会发生泛化，使许多其他状态的估计价值也随之改变。
以这种方式学习模仿输入--输出样例的机器学习方法称为\textbf{监督学习}方法；当输出是像 $u$ 这样的数值时，这个过程通常称为\textbf{函数近似}。
函数近似方法期望接收它们试图近似的函数所应具有的输入--输出行为样例。
我们把每次更新中的 $s\mapsto g$ 作为一个训练样例传递给这些方法，就可以把它们用于价值预测。
然后，我们把它们产生的近似函数解释为估计价值函数。

以这种方式把每次更新看作一个常规训练样例，使我们能够把大量已有函数近似方法中的任一种用于价值预测。
原则上，我们可以使用任何从样例进行监督学习的方法，包括人工神经网络、决策树以及各种多元回归。
然而，并非所有函数近似方法都同样适合用于强化学习。
最复杂的人工神经网络方法和统计方法通常都假设有一个静态训练集，并在其上进行多次遍历。
但是在强化学习中，重要的是学习能够在线发生，也就是在智能体与其环境或环境模型交互时进行。
要做到这一点，需要能够高效地从增量获得的数据中学习的方法。
此外，强化学习通常要求函数近似方法能够处理非平稳目标函数（即会随时间变化的目标函数）。
例如，在基于 GPI（广义策略迭代）的控制方法中，我们常常在 $\pi$ 改变的同时试图学习 $q_\pi$。
即使策略保持不变，如果训练样例的目标值由自举方法（DP 和 TD 学习）生成，那么它们也是非平稳的。
不能轻易处理这种非平稳性的方法，不太适合强化学习。

\subsection*{9.2\quad 预测目标（$\overline{\mathrm{VE}}$）}

到目前为止，我们还没有为预测指定一个明确目标。
在表格型情形中，不需要连续度量预测质量，因为学到的价值函数可以准确等于真实价值函数。
此外，每个状态处的学习价值彼此解耦，一个状态上的更新不会影响其他状态。
但是在真正的近似情形中，一个状态上的更新会影响许多其他状态，而且不可能让所有状态的价值都完全正确。
按照假设，状态数远多于权重数，因此让一个状态的估计更准确，通常必然意味着让另一些状态的估计更不准确。
于是，我们必须说明自己最关心哪些状态。
我们必须指定一个状态分布 $\mu(s)\ge 0$、$\sum_s\mu(s)=1$，用它表示我们对每个状态 $s$ 中误差的关心程度。
所谓状态 $s$ 中的误差，是指近似价值 $\hat{v}(s,\mathbf{w})$ 与真实价值 $v_\pi(s)$ 之差的平方。
用 $\mu$ 在状态空间上加权后，我们得到一个自然的目标函数，即均方价值误差，记为 $\overline{\mathrm{VE}}$：
\begin{equation}
\overline{\mathrm{VE}}(\mathbf{w})\doteq
\sum_{s\in\mathcal{S}}\mu(s)\left[v_\pi(s)-\hat{v}(s,\mathbf{w})\right]^2.
\tag{9.1}
\end{equation}

这个度量的平方根，即 $\overline{\mathrm{VE}}$ 的平方根，给出了近似价值与真实价值相差多少的一个粗略度量，并且常用于作图。
通常，$\mu(s)$ 取为在状态 $s$ 中花费的时间比例。
在同策略训练下，这称为\textbf{同策略分布}；本章将完全关注这种情形。
在连续任务中，同策略分布就是策略 $\pi$ 下的平稳分布。

\begin{center}
\fbox{\begin{minipage}{0.9\textwidth}
\noindent\textbf{分幕任务中的同策略分布}

\medskip
在分幕任务中，同策略分布略有不同，因为它依赖于每一幕初始状态如何选择。
令 $h(s)$ 表示一幕从每个状态 $s$ 开始的概率，令 $\eta(s)$ 表示在一幕中平均花在状态 $s$ 中的时间步数。
如果幕从 $s$ 开始，或者如果从某个花费时间的前一状态 $\bar{s}$ 转移到 $s$，那么时间就花在状态 $s$ 中：
\begin{equation}
\eta(s)=h(s)+\sum_{\bar{s}}\eta(\bar{s})\sum_a\pi(a|\bar{s})p(s|\bar{s},a),
\quad \text{对所有 }s\in\mathcal{S}.
\tag{9.2}
\end{equation}
这个方程组可以解出期望访问次数 $\eta(s)$。
于是，同策略分布就是每个状态中花费时间比例的归一化形式，使其和为一：
\begin{equation}
\mu(s)=\frac{\eta(s)}{\sum_{s'}\eta(s')},
\quad \text{对所有 }s\in\mathcal{S}.
\tag{9.3}
\end{equation}
这是没有折扣时的自然选择。
如果存在折扣（$\gamma<1$），则应当把它看作一种终止形式；这只需在 (9.2) 的第二项中包含一个因子 $\gamma$ 即可。
\end{minipage}}
\end{center}

连续与分幕这两种情形的行为相似；但是在采用近似时，正式分析中必须分别处理它们，这一点我们会在本书这一部分反复看到。
至此，学习目标的规定就完成了。

不过，$\overline{\mathrm{VE}}$ 是否是强化学习中正确的性能目标，并不完全清楚。
请记住，我们的最终目的，也就是我们学习价值函数的原因，是找到更好的策略。
为此目的而言，最好的价值函数并不一定就是使 $\overline{\mathrm{VE}}$ 最小的那个。
尽管如此，至少还不清楚价值预测还有什么更有用的替代目标。
目前，我们将集中讨论 $\overline{\mathrm{VE}}$。

就 $\overline{\mathrm{VE}}$ 而言，一个理想目标是找到全局最优，即某个权重向量 $\mathbf{w}^*$，使得对所有可能的 $\mathbf{w}$ 都有
\[
\overline{\mathrm{VE}}(\mathbf{w}^*)\le \overline{\mathrm{VE}}(\mathbf{w}).
\]
对于线性函数近似器这样的简单函数近似器，达到这个目标有时是可能的；但对于人工神经网络和决策树这样的复杂函数近似器，这很少可能。
退一步说，复杂函数近似器可以转而寻求收敛到局部最优，即某个权重向量 $\mathbf{w}^*$，使得对 $\mathbf{w}^*$ 某个邻域中的所有 $\mathbf{w}$ 都有
\[
\overline{\mathrm{VE}}(\mathbf{w}^*)\le \overline{\mathrm{VE}}(\mathbf{w}).
\]
虽然这样的保证只能稍微令人安心，但它通常已经是非线性函数近似器所能给出的最好说法，而且往往也足够。
不过，在强化学习中许多值得关注的情形下，并不能保证收敛到某个最优，甚至不能保证收敛到距某个最优有界的范围内。
有些方法实际上可能发散，使它们的 $\overline{\mathrm{VE}}$ 在极限中趋于无穷。

在前两节中，我们勾勒了一个框架，用于把大量价值预测强化学习方法与大量函数近似方法结合起来：前者的更新为后者生成训练样例。
我们还描述了一个 $\overline{\mathrm{VE}}$ 性能度量，这些方法可以把它作为希望最小化的目标。
可能的函数近似方法范围太广，无法一一覆盖，而且无论如何，我们对其中大多数方法了解得太少，无法给出可靠评价或建议。
因此，我们只能考察少数可能性。
在本章余下部分，我们关注基于梯度原理的函数近似方法，尤其是线性梯度下降方法。
我们关注这些方法，一方面是因为我们认为它们特别有前景，并且能够揭示关键理论问题；另一方面也是因为它们简单，而我们的篇幅有限。

\subsection*{9.3\quad 随机梯度方法与半梯度方法}

现在，我们详细发展一类用于价值预测中函数近似的学习方法，也就是基于随机梯度下降（stochastic gradient descent, SGD）的方法。
SGD 方法是所有函数近似方法中使用最广泛的一类，并且特别适合在线强化学习。

在梯度下降方法中，权重向量是一个具有固定个数实值分量的列向量，
\[
\mathbf{w}\doteq (w_1,w_2,\ldots,w_d)^\top,
\]
而近似价值函数 $\hat{v}(s,\mathbf{w})$ 对所有 $s\in\mathcal{S}$ 都是 $\mathbf{w}$ 的可微函数。\footnote{$\top$ 表示转置；这里需要用它把正文中的水平行向量变成竖直列向量。在本书中，除非显式横向写出或转置，向量通常都取为列向量。}
我们将在一系列离散时间步 $t=0,1,2,3,\ldots$ 上更新 $\mathbf{w}$，因此需要用记号 $\mathbf{w}_t$ 表示每一步的权重向量。
现在，假设在每一步上，我们观察到一个新的样例 $S_t\mapsto v_\pi(S_t)$，它由一个（可能随机选择的）状态 $S_t$ 及其在该策略下的真实价值组成。
这些状态可能是与环境交互时相继出现的状态，但暂时我们不作这样的假设。
即使给定了每个 $S_t$ 的精确正确价值 $v_\pi(S_t)$，仍然存在一个困难问题，因为我们的函数近似器资源有限，因而分辨率有限。
特别地，通常不存在某个 $\mathbf{w}$ 能够让所有状态，甚至所有样例，都完全正确。
此外，我们还必须泛化到所有未在样例中出现过的其他状态。

我们假设样例中的状态按照同一个分布 $\mu$ 出现，而我们试图最小化的正是由 (9.1) 给出的 $\overline{\mathrm{VE}}$。
在这种情况下，一个好策略是试图最小化已观察样例上的误差。
\textbf{随机梯度下降}（SGD）方法通过在每个样例之后，把权重向量沿着最能降低该样例误差的方向小幅调整来做到这一点：
\begin{align}
\mathbf{w}_{t+1}
&\doteq \mathbf{w}_t-\frac{1}{2}\alpha\nabla\!\left[v_\pi(S_t)-\hat{v}(S_t,\mathbf{w}_t)\right]^2 \tag{9.4}\\
&= \mathbf{w}_t+\alpha\left[v_\pi(S_t)-\hat{v}(S_t,\mathbf{w}_t)\right]\nabla\hat{v}(S_t,\mathbf{w}_t),
\tag{9.5}
\end{align}
其中 $\alpha$ 是正的步长参数；对于任何作为向量（这里是 $\mathbf{w}$）的函数的标量表达式 $f(\mathbf{w})$，$\nabla f(\mathbf{w})$ 表示该表达式对向量各分量的偏导数组成的列向量：
\begin{equation}
\nabla f(\mathbf{w})\doteq
\left(
\frac{\partial f(\mathbf{w})}{\partial w_1},
\frac{\partial f(\mathbf{w})}{\partial w_2},
\ldots,
\frac{\partial f(\mathbf{w})}{\partial w_d}
\right)^\top.
\tag{9.6}
\end{equation}
这个导数向量就是 $f$ 关于 $\mathbf{w}$ 的\textbf{梯度}。
SGD 方法之所以是“梯度下降”方法，是因为 $\mathbf{w}_t$ 中的总体步长与样例平方误差 (9.4) 的负梯度成比例。
这就是误差下降最快的方向。
当更新像这里一样只基于单个样例进行，而这个样例可能是随机选出的时，梯度下降方法称为“随机”的。
在许多样例上不断迈出小步，其总体效果就是最小化像 $\overline{\mathrm{VE}}$ 这样的平均性能度量。

SGD 为什么只沿梯度方向迈出一小步，乍看之下也许并不明显。
难道不能沿这个方向一直走到底，并完全消除该样例上的误差吗？
在许多情形下这可以做到，但通常并不可取。
请记住，我们既不寻求、也不期望找到一个对所有状态误差为零的价值函数，而只是寻找一个平衡不同状态误差的近似。
如果一步就完全修正每个样例，就不会得到这样的平衡。
事实上，SGD 方法的收敛结果假设 $\alpha$ 会随时间减小。
如果它以满足标准随机近似条件 (2.7) 的方式减小，那么 SGD 方法 (9.5) 保证收敛到局部最优。

现在转向这样一种情形：第 $t$ 个训练样例 $S_t\mapsto U_t$ 的目标输出（这里记为 $U_t\in\mathbb{R}$）不是真实价值 $v_\pi(S_t)$，而只是对它的某种可能随机的近似。
例如，$U_t$ 可以是 $v_\pi(S_t)$ 的一个带噪版本，也可以是上一节提到的某个使用 $\hat{v}$ 的自举目标。
在这些情形中，我们无法执行精确更新 (9.5)，因为 $v_\pi(S_t)$ 未知；但我们可以用 $U_t$ 代替 $v_\pi(S_t)$ 来近似它。
这就得到如下用于状态价值预测的一般 SGD 方法：
\begin{equation}
\mathbf{w}_{t+1}\doteq
\mathbf{w}_t+\alpha\left[U_t-\hat{v}(S_t,\mathbf{w}_t)\right]\nabla\hat{v}(S_t,\mathbf{w}_t).
\tag{9.7}
\end{equation}
如果 $U_t$ 是无偏估计，也就是说，对每个 $t$，若 $\mathbb{E}[U_t|S_t=s]=v_\pi(S_t)$，那么在通常的、$\alpha$ 递减的随机近似条件 (2.7) 下，$\mathbf{w}_t$ 保证收敛到局部最优。

例如，假设样例中的状态就是使用策略 $\pi$ 与环境交互（或模拟交互）所生成的状态。
由于一个状态的真实价值是其后续回报的期望，蒙特卡洛目标 $U_t\doteq G_t$ 按定义就是 $v_\pi(S_t)$ 的无偏估计。
在这种选择下，一般 SGD 方法 (9.7) 会收敛到 $v_\pi(S_t)$ 的局部最优近似。
因此，蒙特卡洛状态价值预测的梯度下降版本保证能够找到一个局部最优解。
完整算法的伪代码如下框所示。

\begin{center}
\fbox{\begin{minipage}{0.9\textwidth}
\noindent\textbf{梯度蒙特卡洛算法，用于估计 $\hat{v}\approx v_\pi$}

\medskip
\noindent 输入：待评估的策略 $\pi$。

\noindent 输入：可微函数 $\hat{v}:\mathcal{S}\times\mathbb{R}^d\to\mathbb{R}$。

\noindent 算法参数：步长 $\alpha>0$。

\noindent 任意初始化价值函数权重 $\mathbf{w}\in\mathbb{R}^d$（例如 $\mathbf{w}=\mathbf{0}$）。

\medskip
\noindent 永远循环（对每一幕）：

\quad 用 $\pi$ 生成一幕 $S_0,A_0,R_1,S_1,A_1,\ldots,R_T,S_T$。

\quad 对幕中的每一步 $t=0,1,\ldots,T-1$ 循环：

\qquad $\mathbf{w}\leftarrow \mathbf{w}+\alpha\left[G_t-\hat{v}(S_t,\mathbf{w})\right]\nabla\hat{v}(S_t,\mathbf{w})$。
\end{minipage}}
\end{center}

如果在 (9.7) 中把 $v_\pi(S_t)$ 的自举估计用作目标 $U_t$，就得不到同样的保证。
诸如 $n$ 步回报 $G_{t:t+n}$，或 DP 目标
\[
\sum_{a,s',r}\pi(a|S_t)p(s',r|S_t,a)\left[r+\gamma\hat{v}(s',\mathbf{w}_t)\right]
\]
这样的自举目标，都依赖于当前权重向量 $\mathbf{w}_t$ 的值；这意味着它们会有偏，而且不会产生真正的梯度下降方法。
一种理解方式是，从 (9.4) 到 (9.5) 的关键一步依赖于目标与 $\mathbf{w}_t$ 无关。
如果用自举估计代替 $v_\pi(S_t)$，这一步就不再成立。
自举方法实际上并不是真正梯度下降的实例（Barnard, 1993）。
它们考虑改变权重向量 $\mathbf{w}_t$ 对估计值的影响，却忽略它对目标的影响。
它们只包含梯度的一部分，因此我们称其为\textbf{半梯度方法}。

虽然半梯度（自举）方法不像梯度方法那样稳健地收敛，但在重要情形中，例如下一节讨论的线性情形，它们确实可靠收敛。
此外，它们具有重要优势，因而常常明显更受偏好。
原因之一是它们通常能显著加快学习，正如我们在第 6 章和第 7 章已经看到的那样。
另一个原因是，它们使学习可以连续、在线进行，而不必等待一幕结束。
这使它们可以用于连续问题，并提供计算优势。
一种典型的半梯度方法是半梯度 TD(0)，它使用
\[
U_t\doteq R_{t+1}+\gamma\hat{v}(S_{t+1},\mathbf{w})
\]
作为目标。
该方法的完整伪代码如下框所示。

\begin{center}
\fbox{\begin{minipage}{0.9\textwidth}
\noindent\textbf{半梯度 TD(0)，用于估计 $\hat{v}\approx v_\pi$}

\medskip
\noindent 输入：待评估的策略 $\pi$。

\noindent 输入：可微函数 $\hat{v}:\mathcal{S}^+\times\mathbb{R}^d\to\mathbb{R}$，且 $\hat{v}(\mathrm{terminal},\cdot)=0$。

\noindent 算法参数：步长 $\alpha>0$。

\noindent 任意初始化价值函数权重 $\mathbf{w}\in\mathbb{R}^d$（例如 $\mathbf{w}=\mathbf{0}$）。

\medskip
\noindent 对每一幕循环：

\quad 初始化 $S$。

\quad 对幕中的每一步循环：

\qquad 选择 $A\sim \pi(\cdot|S)$。

\qquad 采取动作 $A$，观察 $R,S'$。

\qquad $\mathbf{w}\leftarrow \mathbf{w}+\alpha\left[R+\gamma\hat{v}(S',\mathbf{w})-\hat{v}(S,\mathbf{w})\right]\nabla\hat{v}(S,\mathbf{w})$。

\qquad $S\leftarrow S'$。

\quad 直到 $S$ 为终止状态。
\end{minipage}}
\end{center}

\textbf{状态聚合}是一种简单的泛化函数近似形式；在其中，状态被分组，每一组有一个估计价值（即权重向量 $\mathbf{w}$ 的一个分量）。
一个状态的价值估计为其所在组的那个分量；当该状态被更新时，也只更新这个分量。
状态聚合是 SGD (9.7) 的一个特例，其中梯度 $\nabla\hat{v}(S_t,\mathbf{w}_t)$ 对 $S_t$ 所在组的分量为 $1$，对其他分量为 $0$。

\noindent\textbf{例 9.1：1000 状态随机游走上的状态聚合}
考虑随机游走任务的 1000 状态版本（第 125 页例 6.2 和第 144 页例 7.1）。
状态从左到右编号为 $1$ 到 $1000$，所有幕都从中心附近的状态 $500$ 开始。
状态转移从当前状态到其左侧 100 个邻近状态之一，或到其右侧 100 个邻近状态之一，所有选择概率相等。
当然，如果当前状态靠近边界，那么这一侧的邻近状态可能少于 100 个。
在这种情况下，本应分配给那些缺失邻近状态的全部概率，都转入在该侧终止的概率（因此，状态 $1$ 有 $0.5$ 的概率在左侧终止，而状态 $950$ 有 $0.25$ 的概率在右侧终止）。
照常，左侧终止产生奖励 $-1$，右侧终止产生奖励 $+1$。
所有其他转移的奖励均为 $0$。
我们将在本节中把这个任务作为贯穿示例。

图 9.1 显示了该任务的真实价值函数 $v_\pi$。
它几乎是一条直线，只是在两端最后 100 个状态附近略微向水平弯曲。
图中还显示了梯度蒙特卡洛算法在状态聚合下、经过 100000 幕并使用步长 $\alpha=2\times10^{-5}$ 后学到的最终近似价值函数。
对于状态聚合，1000 个状态被划分为 10 组，每组 100 个状态（也就是说，状态 1--100 为一组，状态 101--200 为另一组，依此类推）。
图中显示的阶梯状效应是状态聚合的典型表现：在每一组内部，近似价值是常数，而从一组到下一组时会突然改变。
这些近似价值接近 $\overline{\mathrm{VE}}$ (9.1) 的全局最小值。

\begin{figure}[htbp]
\centering
\includegraphics[width=0.78\textwidth]{tex/强化学习/figures/math-10/figure-9-1.png}
\caption{在 1000 状态随机游走任务上，使用梯度蒙特卡洛算法（第 202 页）和状态聚合进行函数近似。}
\end{figure}

近似价值的一些细节最好参照该任务的状态分布 $\mu$ 来理解；该分布显示在图的下部，并使用右侧刻度。
中心处的状态 500 是每一幕的起始状态，但此后很少再次被访问。
平均而言，大约 $1.37\%$ 的时间步花在起始状态中。
从起始状态一步可达的状态是第二常访问的状态，每个这样的状态大约占 $0.17\%$ 的时间步。
从那里开始，$\mu$ 几乎线性下降，在两端状态 1 和 1000 处达到约 $0.0147\%$。
该分布最明显的影响出现在最左侧的组，其价值明显高于该组内真实价值的无权平均；也出现在最右侧的组，其价值明显更低。
这是因为这些区域中的状态在 $\mu$ 下的加权最不对称。
例如，在最左侧组中，状态 100 的权重大约比状态 1 强 3 倍以上。
因此，该组的估计偏向状态 100 的真实价值，而状态 100 的真实价值高于状态 1 的真实价值。

\subsection*{9.4\quad 线性方法}

函数近似中最重要的特殊情形之一，是近似函数 $\hat{v}(\cdot,\mathbf{w})$ 是权重向量 $\mathbf{w}$ 的线性函数。
对应每个状态 $s$，有一个实值向量
\[
\mathbf{x}(s)\doteq (x_1(s),x_2(s),\ldots,x_d(s))^\top,
\]
它的分量数与 $\mathbf{w}$ 相同。
线性方法通过
$\mathbf{w}$ 与 $\mathbf{x}(s)$ 的内积来近似状态价值函数：
\begin{equation}
\hat{v}(s,\mathbf{w})\doteq
\mathbf{w}^\top\mathbf{x}(s)\doteq \sum_{i=1}^d w_i x_i(s).
\tag{9.8}
\end{equation}
在这种情况下，称近似价值函数是\textbf{关于权重线性的}，或简称为\textbf{线性的}。

向量 $\mathbf{x}(s)$ 称为表示状态 $s$ 的\textbf{特征向量}。
$\mathbf{x}(s)$ 的每个分量 $x_i(s)$ 都是某个函数 $x_i:\mathcal{S}\to\mathbb{R}$ 的取值。
我们把这些函数中的每一个整体看作一个\textbf{特征}，并把它在状态 $s$ 处的值称为 $s$ 的一个特征。
对于线性方法，特征就是\textbf{基函数}，因为它们为近似函数集合形成一个线性基。
构造 $d$ 维特征向量来表示状态，就等价于选择一组 $d$ 个基函数。
特征可以用许多不同方式定义；接下来的几节将介绍几种可能性。

在线性函数近似中使用 SGD 更新是很自然的。
在这种情况下，近似价值函数关于 $\mathbf{w}$ 的梯度为
\[
\nabla\hat{v}(s,\mathbf{w})=\mathbf{x}(s).
\]
因此，在线性情形中，一般 SGD 更新 (9.7) 化为特别简单的形式：
\[
\mathbf{w}_{t+1}\doteq
\mathbf{w}_t+\alpha\left[U_t-\hat{v}(S_t,\mathbf{w}_t)\right]\mathbf{x}(S_t).
\]
正因为它如此简单，线性 SGD 情形是最有利于数学分析的情形之一。
几乎所有关于各类学习系统的有用收敛结果，都是针对线性（或更简单的）函数近似方法得到的。
特别地，在线性情形中只有一个最优（或者在退化情形中，有一组同样好的最优），因此任何保证收敛到局部最优或其附近的方法，都会自动保证收敛到全局最优或其附近。
例如，上一节给出的梯度蒙特卡洛算法，在线性函数近似下，如果 $\alpha$ 按通常条件随时间减小，就会收敛到 $\overline{\mathrm{VE}}$ 的全局最优。

上一节给出的半梯度 TD(0) 算法在线性函数近似下也会收敛，但这并不能由 SGD 的一般结果推出；这里需要一个单独的定理。
它收敛到的权重向量也不是全局最优，而是接近局部最优的一个点。
更细致地考察这个重要情形是有用的，尤其是在连续任务中。
每个时刻 $t$ 的更新为
\begin{align}
\mathbf{w}_{t+1}
&\doteq \mathbf{w}_t+\alpha\left(R_{t+1}+\gamma\mathbf{w}_t^\top\mathbf{x}_{t+1}-\mathbf{w}_t^\top\mathbf{x}_t\right)\mathbf{x}_t \tag{9.9}\\
&= \mathbf{w}_t+\alpha\left(R_{t+1}\mathbf{x}_t-\mathbf{x}_t(\mathbf{x}_t-\gamma\mathbf{x}_{t+1})^\top\mathbf{w}_t\right),
\end{align}
其中我们使用了记号简写 $\mathbf{x}_t=\mathbf{x}(S_t)$。
一旦系统达到稳态，对于任意给定的 $\mathbf{w}_t$，期望的下一权重向量可写为
\begin{equation}
\mathbb{E}[\mathbf{w}_{t+1}|\mathbf{w}_t]
=\mathbf{w}_t+\alpha(\mathbf{b}-\mathbf{A}\mathbf{w}_t),
\tag{9.10}
\end{equation}
其中
\begin{equation}
\mathbf{b}\doteq \mathbb{E}[R_{t+1}\mathbf{x}_t]\in\mathbb{R}^d,
\qquad
\mathbf{A}\doteq \mathbb{E}\!\left[\mathbf{x}_t(\mathbf{x}_t-\gamma\mathbf{x}_{t+1})^\top\right]\in\mathbb{R}^d\times\mathbb{R}^d.
\tag{9.11}
\end{equation}
由 (9.10) 可清楚看出，如果系统收敛，那么它必然收敛到权重向量 $\mathbf{w}_{\mathrm{TD}}$，并在该处满足
\begin{align}
\mathbf{b}-\mathbf{A}\mathbf{w}_{\mathrm{TD}}&=0 \notag\\
\Rightarrow\quad \mathbf{b}&=\mathbf{A}\mathbf{w}_{\mathrm{TD}} \notag\\
\Rightarrow\quad \mathbf{w}_{\mathrm{TD}}&\doteq \mathbf{A}^{-1}\mathbf{b}.
\tag{9.12}
\end{align}
这个量称为 \textbf{TD 不动点}。
事实上，线性半梯度 TD(0) 会收敛到这一点。
证明其收敛性以及上面逆矩阵存在性的一部分理论见下框。

\begin{center}
\fbox{\begin{minipage}{0.9\textwidth}
\noindent\textbf{线性 TD(0) 收敛性的证明}

\medskip
什么性质能保证线性 TD(0) 算法 (9.9) 收敛？
把 (9.10) 改写为如下形式，可以得到一些直觉：
\begin{equation}
\mathbb{E}[\mathbf{w}_{t+1}|\mathbf{w}_t]
=(\mathbf{I}-\alpha\mathbf{A})\mathbf{w}_t+\alpha\mathbf{b}.
\tag{9.13}
\end{equation}
注意，矩阵 $\mathbf{A}$ 乘的是权重向量 $\mathbf{w}_t$，而不是 $\mathbf{b}$；只有 $\mathbf{A}$ 对收敛性重要。
为了形成直觉，考虑 $\mathbf{A}$ 是对角矩阵的特殊情形。
如果某个对角元素为负，那么 $\mathbf{I}-\alpha\mathbf{A}$ 中相应的对角元素将大于 $1$，$\mathbf{w}_t$ 的相应分量会被放大；如果这种情况持续，就会导致发散。
另一方面，如果 $\mathbf{A}$ 的所有对角元素均为正，那么可以选择 $\alpha$ 小于其中最大元素的倒数，使得 $\mathbf{I}-\alpha\mathbf{A}$ 是一个所有对角元素都位于 $0$ 与 $1$ 之间的对角矩阵。
在这种情况下，更新的第一项倾向于收缩 $\mathbf{w}_t$，稳定性得到保证。
一般而言，只要 $\mathbf{A}$ 是正定的，$\mathbf{w}_t$ 就会向零收缩；正定是指对任意实向量 $\mathbf{y}\ne 0$，都有 $\mathbf{y}^\top\mathbf{A}\mathbf{y}>0$。
正定性还保证逆矩阵 $\mathbf{A}^{-1}$ 存在。

对于连续情形且 $\gamma<1$ 的线性 TD(0)，矩阵 $\mathbf{A}$ (9.11) 可以写为
\begin{align*}
\mathbf{A}
&=\sum_s\mu(s)\sum_a\pi(a|s)\sum_{r,s'}p(r,s'|s,a)\mathbf{x}(s)\left(\mathbf{x}(s)-\gamma\mathbf{x}(s')\right)^\top\\
&=\sum_s\mu(s)\sum_{s'}p(s'|s)\mathbf{x}(s)\left(\mathbf{x}(s)-\gamma\mathbf{x}(s')\right)^\top\\
&=\sum_s\mu(s)\mathbf{x}(s)
\left(\mathbf{x}(s)-\gamma\sum_{s'}p(s'|s)\mathbf{x}(s')\right)^\top\\
&=\mathbf{X}^\top\mathbf{D}(\mathbf{I}-\gamma\mathbf{P})\mathbf{X},
\end{align*}
其中 $\mu(s)$ 是策略 $\pi$ 下的平稳分布，$p(s'|s)$ 是在策略 $\pi$ 下从 $s$ 转移到 $s'$ 的概率，$\mathbf{P}$ 是由这些概率组成的 $|\mathcal{S}|\times|\mathcal{S}|$ 矩阵，$\mathbf{D}$ 是对角线上为 $\mu(s)$ 的 $|\mathcal{S}|\times|\mathcal{S}|$ 对角矩阵，$\mathbf{X}$ 是以 $\mathbf{x}(s)$ 为各行的 $|\mathcal{S}|\times d$ 矩阵。
由此可见，内层矩阵 $\mathbf{D}(\mathbf{I}-\gamma\mathbf{P})$ 是决定 $\mathbf{A}$ 正定性的关键。

对于这种形式的关键矩阵，如果它的所有列和都是非负数，则正定性得到保证。
这一点由 Sutton（1988，第 27 页）基于两个已发表定理证明。
第一个定理说，任意矩阵 $\mathbf{M}$ 是正定的，当且仅当对称矩阵 $\mathbf{S}=\mathbf{M}+\mathbf{M}^\top$ 是正定的（Sutton 1988，附录）。
第二个定理说，任意对称实矩阵 $\mathbf{S}$，若其所有对角元素都为正，并且大于相应非对角元素绝对值之和，则它是正定的（Varga 1962，第 23 页）。
对于我们的关键矩阵 $\mathbf{D}(\mathbf{I}-\gamma\mathbf{P})$，对角元素为正而非对角元素为负，因此只需证明每一行和加上相应列和为正即可。
行和全为正，因为 $\mathbf{P}$ 是随机矩阵且 $\gamma<1$。
于是只剩下证明列和非负。
注意，任意矩阵 $\mathbf{M}$ 的列和组成的行向量可以写作 $\mathbf{1}^\top\mathbf{M}$，其中 $\mathbf{1}$ 是所有分量都等于 $1$ 的列向量。
令 $\boldsymbol{\mu}$ 表示由 $\mu(s)$ 组成的 $|\mathcal{S}|$ 维向量；由于 $\mu$ 是平稳分布，有 $\boldsymbol{\mu}=\mathbf{P}^\top\boldsymbol{\mu}$。
于是，我们关键矩阵的列和为
\begin{align*}
\mathbf{1}^\top\mathbf{D}(\mathbf{I}-\gamma\mathbf{P})
&=\boldsymbol{\mu}^\top(\mathbf{I}-\gamma\mathbf{P})\\
&=\boldsymbol{\mu}^\top-\gamma\boldsymbol{\mu}^\top\mathbf{P}\\
&=\boldsymbol{\mu}^\top-\gamma\boldsymbol{\mu}^\top
\qquad\text{（因为 $\boldsymbol{\mu}$ 是平稳分布）}\\
&=(1-\gamma)\boldsymbol{\mu}^\top,
\end{align*}
其中所有分量都是正的。
因此，关键矩阵及其 $\mathbf{A}$ 矩阵都是正定的，同策略 TD(0) 是稳定的。
（要证明以概率 1 收敛，还需要关于 $\alpha$ 随时间减小的附加条件和计划。）
\end{minipage}}
\end{center}

在 TD 不动点处，还可以证明（在连续情形中）$\overline{\mathrm{VE}}$ 位于最低可能误差的一个有界放大范围内：
\begin{equation}
\overline{\mathrm{VE}}(\mathbf{w}_{\mathrm{TD}})
\le \frac{1}{1-\gamma}\min_{\mathbf{w}}\overline{\mathrm{VE}}(\mathbf{w}).
\tag{9.14}
\end{equation}
也就是说，TD 方法的渐近误差至多是最小可能误差的 $\frac{1}{1-\gamma}$ 倍；这个最小可能误差是蒙特卡洛方法在极限中达到的误差。
因为 $\gamma$ 往往接近 $1$，这个放大因子可能相当大，所以 TD 方法的渐近性能可能有实质损失。
另一方面，回忆一下，与蒙特卡洛方法相比，TD 方法的方差往往大幅降低，因此学习更快，正如我们在第 6 章和第 7 章所见。
哪种方法最好，取决于近似和问题的性质，也取决于学习持续多久。

与 (9.14) 类似的界也适用于其他同策略自举方法。
例如，若线性半梯度 DP（公式 9.7，并取
\[
U_t\doteq \sum_a\pi(a|S_t)\sum_{s',r}p(s',r|S_t,a)\left[r+\gamma\hat{v}(s',\mathbf{w}_t)\right]
\]
）按照同策略分布进行更新，也会收敛到 TD 不动点。
一步半梯度动作价值方法，例如下一章介绍的半梯度 Sarsa(0)，会收敛到一个类似的不动点，并有一个类似的界。
对于分幕任务，还存在一个略有不同但相关的界（见 Bertsekas and Tsitsiklis, 1996）。
这里还省略了一些关于奖励、特征以及步长参数递减的技术条件。
完整细节可见原始论文（Tsitsiklis and Van Roy, 1997）。

这些收敛结果的关键在于，状态按照同策略分布进行更新。
对于其他更新分布，使用函数近似的自举方法实际上可能发散到无穷。
相关例子以及可能解法的讨论见第 11 章。

\noindent\textbf{例 9.2：1000 状态随机游走上的自举}
状态聚合是线性函数近似的一个特例，因此让我们回到 1000 状态随机游走，以说明本章作出的一些观察。
图 9.2 左图显示半梯度 TD(0) 算法（第 203 页）在使用与例 9.1 相同的状态聚合时学到的最终价值函数。
可以看到，近似渐近的 TD 近似确实比图 9.1 所示的蒙特卡洛近似离真实价值更远。

尽管如此，TD 方法在学习速度上仍然保有很大的潜在优势，并且可以推广蒙特卡洛方法，正如我们在第 7 章用 $n$ 步 TD 方法已经充分考察过的那样。
图 9.2 右图显示了在 1000 状态随机游走上使用状态聚合的 $n$ 步半梯度 TD 方法的结果；它们与我们先前用表格方法和 19 状态随机游走（图 7.2）得到的结果惊人相似。
为了得到这种定量上相近的结果，我们把状态聚合改成 20 组，每组 50 个状态。
这样，这 20 组就与表格问题中的 19 个状态在定量上接近。
特别地，回忆一下，状态转移最多向左或向右 100 个状态。
于是，典型转移会向左或向右移动 50 个状态，这与 19 状态表格系统中的单状态转移在定量上类似。
为了使匹配完整，我们这里使用相同性能度量，即前 10 幕中所有状态 RMS 误差的无权平均，而不是在使用函数近似时通常更合适的 $\overline{\mathrm{VE}}$ 目标。

\begin{figure}[htbp]
\centering
\includegraphics[width=0.9\textwidth]{tex/强化学习/figures/math-10/figure-9-2.png}
\caption{在 1000 状态随机游走任务上，用状态聚合进行自举。左：半梯度 TD 的渐近价值比图 9.1 中的渐近蒙特卡洛价值更差。右：带状态聚合的 $n$ 步方法的性能与表格表示下的性能惊人相似（参见图 7.2）。这些数据是 100 次运行的平均结果。}
\end{figure}

上例中使用的半梯度 $n$ 步 TD 算法，是第 7 章给出的表格型 $n$ 步 TD 算法向半梯度函数近似的自然扩展。
伪代码如下框所示。

\begin{center}
\fbox{\begin{minipage}{0.92\textwidth}
\noindent\textbf{$n$ 步半梯度 TD，用于估计 $\hat{v}\approx v_\pi$}

\medskip
\noindent 输入：待评估的策略 $\pi$。

\noindent 输入：可微函数 $\hat{v}:\mathcal{S}^+\times\mathbb{R}^d\to\mathbb{R}$，且 $\hat{v}(\mathrm{terminal},\cdot)=0$。

\noindent 算法参数：步长 $\alpha>0$，正整数 $n$。

\noindent 任意初始化价值函数权重 $\mathbf{w}$（例如 $\mathbf{w}=\mathbf{0}$）。

\noindent 所有存储和访问操作（$S_t$ 与 $R_t$）都可以让索引对 $n+1$ 取模。

\medskip
\noindent 对每一幕循环：

\quad 初始化并存储 $S_0\ne \mathrm{terminal}$。

\quad $T\leftarrow\infty$。

\quad 对 $t=0,1,2,\ldots$ 循环：

\qquad 若 $t<T$，则：

\qquad\quad 按照 $\pi(\cdot|S_t)$ 采取一个动作。

\qquad\quad 观察并存储下一奖励为 $R_{t+1}$、下一状态为 $S_{t+1}$。

\qquad\quad 若 $S_{t+1}$ 为终止状态，则 $T\leftarrow t+1$。

\qquad $\tau\leftarrow t-n+1$（$\tau$ 是其状态估计正被更新的时刻）。

\qquad 若 $\tau\ge 0$：

\qquad\quad $G\leftarrow \sum_{i=\tau+1}^{\min(\tau+n,T)}\gamma^{i-\tau-1}R_i$。

\qquad\quad 若 $\tau+n<T$，则 $G\leftarrow G+\gamma^n\hat{v}(S_{\tau+n},\mathbf{w})$（$G_{\tau:\tau+n}$）。

\qquad\quad $\mathbf{w}\leftarrow \mathbf{w}+\alpha\left[G-\hat{v}(S_\tau,\mathbf{w})\right]\nabla\hat{v}(S_\tau,\mathbf{w})$。

\quad 直到 $\tau=T-1$。
\end{minipage}}
\end{center}

这个算法的关键方程类似于 (7.2)，即
\begin{equation}
\mathbf{w}_{t+n}\doteq
\mathbf{w}_{t+n-1}
+\alpha\left[G_{t:t+n}-\hat{v}(S_t,\mathbf{w}_{t+n-1})\right]\nabla\hat{v}(S_t,\mathbf{w}_{t+n-1}),
\quad 0\le t<T,
\tag{9.15}
\end{equation}
其中 $n$ 步回报由 (7.1) 推广为
\begin{equation}
G_{t:t+n}
=R_{t+1}+\gamma R_{t+2}+\cdots+\gamma^{n-1}R_{t+n}
+\gamma^n\hat{v}(S_{t+n},\mathbf{w}_{t+n-1}),
\quad 0\le t\le T-n.
\tag{9.16}
\end{equation}

\noindent\textbf{练习 9.1}
说明本书第一部分介绍的表格型方法是线性函数近似的一个特例。
特征向量会是什么样？

\subsection*{9.5\quad 线性方法的特征构造}

线性方法很有意思，不仅因为它们有收敛保证，也因为在实践中它们在数据和计算两方面都可能非常高效。
是否如此，关键取决于状态如何用特征表示；这正是我们在这一大节中要研究的问题。
为任务选择合适特征，是把先验领域知识加入强化学习系统的一种重要方式。
直观地说，特征应当对应于状态空间中那些适合发生泛化的方面。
例如，如果我们在评价几何对象，可能希望为每一种可能的形状、颜色、大小或功能设置特征。
如果我们在评价移动机器人的状态，那么可能希望为位置、剩余电量程度、最近的声呐读数等设置特征。

线性形式的一个限制是，它不能考虑特征之间的任何交互，例如特征 $i$ 的出现只有在特征 $j$ 不出现时才是好的。
例如，在杆平衡任务（例 3.4）中，高角速度是好是坏取决于角度。
如果角度很大，那么高角速度意味着即将倒下，是坏状态；而如果角度很小，那么高角速度意味着杆正在自行扶正，是好状态。
如果价值函数的特征分别编码角度和角速度，那么一个线性价值函数就无法表示这种关系。
它需要改用，或额外加入，表示这两个底层状态维度组合的特征。
在下面各小节中，我们将考察若干一般做法。

\subsubsection*{9.5.1\quad 多项式}

许多问题中的状态最初都表示为数值，例如杆平衡任务（例 3.4）中的位置和速度，Jack 租车问题（例 4.2）中每个地点的汽车数，或者赌徒问题（例 4.3）中的赌徒资本。
在这些类型的问题中，强化学习的函数近似与熟悉的插值和回归任务有许多共同点。
插值和回归中常用的各种特征族也可以用于强化学习。
多项式构成了用于插值和回归的最简单特征族之一。
虽然这里讨论的基本多项式特征在强化学习中不如其他类型的特征有效，但它们简单而熟悉，因而是很好的引入。

举例来说，假设一个强化学习问题的状态有两个数值维度。
对某个代表性状态 $s$，令它的两个数值为 $s_1\in\mathbb{R}$ 和 $s_2\in\mathbb{R}$。
你可以选择简单地用这两个状态维度表示 $s$，即 $\mathbf{x}(s)=(s_1,s_2)^\top$；但这样就无法考虑这两个维度之间的任何交互。
此外，如果 $s_1$ 和 $s_2$ 都为零，那么近似价值也必须为零。
这两个限制都可以通过改用四维特征向量
\[
\mathbf{x}(s)=(1,s_1,s_2,s_1s_2)^\top
\]
来克服。
开头的 $1$ 特征允许表示原始状态数值上的仿射函数，最后的乘积特征 $s_1s_2$ 则使交互可以被考虑。
或者，也可以使用更高维的特征向量，例如
\[
\mathbf{x}(s)=(1,s_1,s_2,s_1s_2,s_1^2,s_2^2,s_1s_2^2,s_1^2s_2,s_1^2s_2^2)^\top,
\]
以考虑更复杂的交互。
这样的特征向量能够近似任意关于状态数值的二次函数，即使所学习的近似仍然关于权重是线性的。
把这个例子从两个数推广到 $k$ 个数，我们就可以表示问题各状态维度之间高度复杂的交互：

\begin{center}
\fbox{\begin{minipage}{0.9\textwidth}
假设每个状态 $s$ 对应 $k$ 个数 $s_1,s_2,\ldots,s_k$，其中每个 $s_i\in\mathbb{R}$。
对于这个 $k$ 维状态空间，每个 $n$ 阶多项式基特征 $x_i$ 可写为
\begin{equation}
x_i(s)=\prod_{j=1}^k s_j^{c_{i,j}},
\tag{9.17}
\end{equation}
其中每个 $c_{i,j}$ 都是集合 $\{0,1,\ldots,n\}$ 中的整数，且整数 $n\ge 0$。
这些特征构成维度为 $k$ 的 $n$ 阶多项式基，其中包含 $(n+1)^k$ 个不同特征。
\end{minipage}}
\end{center}

更高阶的多项式基可以更准确地近似更复杂的函数。
但是，由于 $n$ 阶多项式基中特征数会随自然状态空间维度 $k$ 指数增长（若 $n>0$），通常有必要选择其中的一个子集用于函数近似。
这可以利用关于待近似函数性质的先验信念完成；一些为多项式回归发展出的自动选择方法，也可以调整后用于处理强化学习的增量性和非平稳性。

\noindent\textbf{练习 9.2}
为什么 (9.17) 对维度 $k$ 定义了 $(n+1)^k$ 个不同特征？

\noindent\textbf{练习 9.3}
什么样的 $n$ 和 $c_{i,j}$ 会产生特征向量
\[
\mathbf{x}(s)=(1,s_1,s_2,s_1s_2,s_1^2,s_2^2,s_1s_2^2,s_1^2s_2,s_1^2s_2^2)^\top?
\]

\subsubsection*{9.5.2\quad 傅里叶基}

另一种线性函数近似方法基于历史悠久的傅里叶级数；傅里叶级数把周期函数表示为不同频率的正弦和余弦基函数（特征）的加权和。
（若对所有 $x$ 以及某个周期 $\tau$ 都有 $f(x)=f(x+\tau)$，则函数 $f$ 是周期的。）
傅里叶级数以及更一般的傅里叶变换广泛用于应用科学，部分原因是如果待近似函数已知，那么基函数权重可以由简单公式给出；进一步说，只要有足够多的基函数，本质上任何函数都可以按需要任意精确地近似。
在函数待近似函数未知的强化学习中，傅里叶基函数之所以值得关注，是因为它们易于使用，并且能在一系列强化学习问题中表现良好。

先考虑一维情形。
一维周期为 $\tau$ 的函数，其通常傅里叶级数表示把该函数表示为正弦与余弦函数的线性组合；这些正弦与余弦函数的周期都能整除 $\tau$（换句话说，其频率是基频 $1/\tau$ 的整数倍）。
但是，如果你关心的是在一个有界区间上定义的非周期函数的近似，那么也可以使用这些傅里叶基特征，并把 $\tau$ 设为该区间长度。
此时，感兴趣的函数只是这些正弦和余弦特征的周期线性组合中的一个周期。

此外，如果把 $\tau$ 设为感兴趣区间长度的两倍，并且只关注半区间 $[0,\tau/2]$ 上的近似，那么只使用余弦特征就可以。
这是可能的，因为任何偶函数，即任何关于原点对称的函数，都可以只用余弦基表示。
因此，只要有足够多的余弦特征，半周期 $[0,\tau/2]$ 上的任何函数都可以按需要任意接近地近似。
（说“任何函数”并不完全正确，因为函数必须在数学上足够良好；这里我们略过这个技术细节。）
另一种选择是只使用正弦特征，因为它们的线性组合总是奇函数，也就是关于原点反对称的函数。
但通常只保留余弦特征更好，因为“半偶”函数往往比“半奇”函数更容易近似，后者经常在原点不连续。
当然，这并不排除在区间 $[0,\tau/2]$ 上同时使用正弦和余弦特征；在某些情况下，这可能有利。

遵循这一逻辑，并令 $\tau=2$，使特征定义在半 $\tau$ 区间 $[0,1]$ 上，一维 $n$ 阶傅里叶余弦基由 $n+1$ 个特征组成：
\[
x_i(s)=\cos(i\pi s),\quad s\in[0,1],
\]
其中 $i=0,\ldots,n$。
图 9.3 显示了一维傅里叶余弦特征 $x_i$，$i=1,2,3,4$；$x_0$ 是常数函数。

\begin{figure}[htbp]
\centering
\includegraphics[width=0.9\textwidth]{tex/强化学习/figures/math-10/figure-9-3.png}
\caption{一维傅里叶余弦基特征 $x_i$，$i=1,2,3,4$，用于近似区间 $[0,1]$ 上的函数。引自 Konidaris 等（2011）。}
\end{figure}

同样的推理也适用于多维情形中的傅里叶余弦级数近似，如下框所述。

\begin{center}
\fbox{\begin{minipage}{0.9\textwidth}
假设每个状态 $s$ 对应一个由 $k$ 个数组成的向量 $\mathbf{s}=(s_1,s_2,\ldots,s_k)^\top$，其中每个 $s_i\in[0,1]$。
$n$ 阶傅里叶余弦基中的第 $i$ 个特征可写为
\begin{equation}
x_i(s)=\cos(\pi \mathbf{s}^\top\mathbf{c}^i),
\tag{9.18}
\end{equation}
其中 $\mathbf{c}^i=(c_1^i,\ldots,c_k^i)^\top$，且对 $j=1,\ldots,k$，$c_j^i\in\{0,\ldots,n\}$，并且 $i=0,\ldots,(n+1)^k$。
这为 $(n+1)^k$ 个可能的整数向量 $\mathbf{c}^i$ 中的每一个定义了一个特征。
内积 $\mathbf{s}^\top\mathbf{c}^i$ 的作用，是给 $s$ 的每个维度分配一个 $\{0,\ldots,n\}$ 中的整数。
和一维情形一样，这个整数决定了特征沿该维度的频率。
这些特征当然可以平移和缩放，以适应特定应用中的有界状态空间。
\end{minipage}}
\end{center}

举例来说，考虑 $k=2$ 的情形，其中 $\mathbf{s}=(s_1,s_2)^\top$，每个 $\mathbf{c}^i=(c_1^i,c_2^i)^\top$。
图 9.4 显示了六个傅里叶余弦特征的选集，每个特征都标出了定义它的向量 $\mathbf{c}^i$（$s_1$ 是水平轴，$\mathbf{c}^i$ 以省略下标 $i$ 的行向量形式显示）。
$\mathbf{c}$ 中的任何零都表示该特征沿相应状态维度是常数。
因此，若 $\mathbf{c}=(0,0)^\top$，则特征在两个维度上都是常数；若 $\mathbf{c}=(c_1,0)^\top$，则特征在第二维上是常数，并在第一维上以取决于 $c_1$ 的频率变化；$\mathbf{c}=(0,c_2)^\top$ 时同理。
当 $\mathbf{c}=(c_1,c_2)^\top$ 且没有 $c_j=0$ 时，特征沿两个维度都变化，并表示两个状态变量之间的交互。
$c_1$ 和 $c_2$ 的取值决定了沿每个维度的频率，而它们的比值给出交互方向。

\begin{figure}[htbp]
\centering
\includegraphics[width=0.9\textwidth]{tex/强化学习/figures/math-10/figure-9-4.png}
\caption{六个二维傅里叶余弦特征的选集，每个特征都用定义它的向量 $\mathbf{c}^i$ 标注（$s_1$ 是水平轴，$\mathbf{c}^i$ 以省略下标 $i$ 的形式显示）。引自 Konidaris 等（2011）。}
\end{figure}

当把傅里叶余弦特征与 (9.7)、半梯度 TD(0) 或半梯度 Sarsa 这样的学习算法结合使用时，为每个特征使用不同步长参数可能会有帮助。
若 $\alpha$ 是基本步长参数，则 Konidaris、Osentoski 和 Thomas（2011）建议把特征 $x_i$ 的步长参数设为
\[
\alpha_i=\frac{\alpha}{\sqrt{(c_1^i)^2+\cdots+(c_k^i)^2}},
\]
但当每个 $c_j^i=0$ 时，取 $\alpha_i=\alpha$。
与若干其他基函数集合相比，包括多项式和径向基函数，傅里叶余弦特征配合 Sarsa 可以产生良好表现。
不过并不令人惊讶的是，傅里叶特征难以处理不连续性，因为除非包含非常高频的基函数，否则很难避免在不连续点周围出现“振铃”。

在 $n$ 阶傅里叶基中，特征数随状态空间维度指数增长；但如果该维度足够小（例如 $k\le5$），就可以选择 $n$，使得所有 $n$ 阶傅里叶特征都可使用。
这使特征选择多少有些自动化。
然而，对于高维状态空间，必须选择这些特征的一个子集。
这可以利用关于待近似函数性质的先验信念来完成；一些自动选择方法也可以改造后用于处理强化学习的增量性和非平稳性。
傅里叶基特征在这方面的一个优点是，通过设置 $\mathbf{c}^i$ 向量以考虑状态变量之间被怀疑存在的交互，并限制 $\mathbf{c}^j$ 向量中的取值，使近似可以滤除被认为是噪声的高频分量，就很容易选择特征。
另一方面，由于傅里叶特征（除少数为零处外）在整个状态空间上非零，它们表示的是状态的全局性质，这可能使寻找表示局部性质的好方法变得困难。

图 9.5 给出了 1000 状态随机游走例子上比较傅里叶基与多项式基的学习曲线。
一般而言，我们不建议把多项式用于在线学习。\footnote{存在比我们讨论的这些多项式更复杂的多项式族，例如不同的正交多项式族；它们可能表现更好，但目前在强化学习中使用它们的经验很少。}

\begin{figure}[htbp]
\centering
\includegraphics[width=0.72\textwidth]{tex/强化学习/figures/math-10/figure-9-5.png}
\caption{1000 状态随机游走中傅里叶基与多项式的比较。图中显示的是梯度蒙特卡洛方法在 5、10 和 20 阶傅里叶基与多项式基下的学习曲线。每种情形的步长参数都经过粗略优化：多项式基取 $\alpha=0.0001$，傅里叶基取 $\alpha=0.00005$。性能度量（纵轴）是均方价值误差 (9.1) 的平方根。}
\end{figure}

\subsubsection*{9.5.3\quad 粗编码}

考虑一个任务，其中状态集合的自然表示是一个连续二维空间。
这类情形的一种表示由与状态空间中圆形区域对应的特征组成，如图 9.6 所示。
如果状态位于某个圆内，那么相应特征取值为 $1$，并称该特征\textbf{出现}；否则该特征取值为 $0$，并称该特征\textbf{缺失}。
这种取 $1$--$0$ 值的特征称为\textbf{二值特征}。
给定一个状态，哪些二值特征出现，就表明该状态位于哪些圆内，从而对其位置作出粗略编码。
用这种彼此重叠的特征来表示状态（虽然它们不必是圆，也不必是二值的）称为\textbf{粗编码}。

\begin{figure}[htbp]
\centering
\includegraphics[width=0.42\textwidth]{tex/强化学习/figures/math-10/figure-9-6.png}
\caption{粗编码。从状态 $s$ 到状态 $s'$ 的泛化取决于它们有多少特征的感受野（这里是圆）重叠。这两个状态有一个共同特征，因此它们之间会有轻微泛化。}
\end{figure}

假设采用线性梯度下降函数近似，考察圆的大小和密度的影响。
每个圆对应一个会受学习影响的权重（$\mathbf{w}$ 的一个分量）。
如果在空间中的某一点训练一个状态，那么所有与该状态相交的圆的权重都会受影响。
因此，由 (9.8) 可知，近似价值函数会在这些圆的并集内所有状态处受影响；一个点与被训练状态“共有”的圆越多，影响就越大，如图 9.6 所示。
如果圆较小，那么泛化距离较短，如图 9.7 左图；如果圆较大，那么泛化距离较长，如图 9.7 中图。
此外，特征的形状将决定泛化的性质。
例如，如果特征不是严格的圆，而是沿某个方向拉长，那么泛化也会受到类似影响，如图 9.7 右图。

\begin{figure}[htbp]
\centering
\includegraphics[width=0.9\textwidth]{tex/强化学习/figures/math-10/figure-9-7.png}
\caption{线性函数近似方法中的泛化由特征感受野的大小与形状决定。这三种情形中的特征数量和密度大致相同。}
\end{figure}

感受野大的特征会带来广泛泛化，但也可能看起来把学到的函数限制为一种粗略近似，无法作出远细于感受野宽度的区分。
幸运的是，事实并非如此。
从一点到另一点的初始泛化确实由感受野大小和形状控制；但是敏锐度，即最终能够达到的最细区分程度，更多由特征总数控制。

\noindent\textbf{例 9.3：粗编码的粗细程度}
本例说明粗编码中感受野大小对学习的影响。
基于粗编码和 (9.7) 的线性函数近似被用来学习一维方波函数（图 9.8 顶部所示）。
该函数的取值被用作目标 $U_t$。
在只有一个维度时，感受野是区间而不是圆。
学习在三种不同大小的区间上重复进行：窄、中等、宽，如图底部所示。
三种情形有相同的特征密度，在待学习函数的范围上大约都有 50 个特征。
训练样例在该范围内均匀随机生成。
步长参数取 $\alpha=\frac{0.2}{n}$，其中 $n$ 是同一时刻出现的特征数。
图 9.8 显示了三种情形在学习过程中的已学函数。
注意，特征宽度在学习早期有强烈影响。
使用宽特征时，泛化倾向于很宽；使用窄特征时，只有每个训练点的近邻被改变，使学到的函数更凹凸不平。
然而，最终学到的函数只受到特征宽度的轻微影响。
感受野形状往往对泛化有强影响，但对渐近解的质量影响很小。

\begin{figure}[htbp]
\centering
\includegraphics[width=0.78\textwidth]{tex/强化学习/figures/math-10/figure-9-8.png}
\caption{特征宽度对初始泛化有强影响（第一行），但对渐近精度影响较弱（最后一行）的例子。}
\end{figure}

\subsubsection*{9.5.4\quad 瓦片编码}

\textbf{瓦片编码}是用于多维连续空间的一种粗编码形式，灵活且计算高效。
它可能是现代顺序数字计算机上最实用的特征表示。

在瓦片编码中，特征的感受野被分组成状态空间的若干划分。
每个这样的划分称为一个\textbf{铺片}，划分中的每个元素称为一个\textbf{瓦片}。
例如，一个二维状态空间最简单的铺片，是图 9.9 左侧所示的均匀网格。
这里的瓦片或感受野是方形，而不是图 9.6 中的圆。
如果只使用这一个铺片，那么白点所表示的状态将由它所在瓦片对应的单个特征来表示；泛化会完全扩展到同一瓦片内的所有状态，而不会扩展到瓦片外的状态。
只用一个铺片时，我们得到的不是粗编码，而只是状态聚合的一种情形。

\begin{figure}[htbp]
\centering
\includegraphics[width=0.9\textwidth]{tex/强化学习/figures/math-10/figure-9-9.png}
\caption{有限二维空间上的多个重叠网格铺片。这些铺片在每个维度上彼此偏移一个均匀量。}
\end{figure}

要获得粗编码的优点，需要重叠的感受野；但按定义，一个划分中的瓦片彼此不重叠。
为了用瓦片编码得到真正的粗编码，需要使用多个铺片，每个铺片相对于其他铺片偏移一个瓦片宽度的分数。
图 9.9 右侧显示了一个包含四个铺片的简单情形。
每个状态，例如白点所示状态，都恰好落在四个铺片中的各一个瓦片内。
这四个瓦片对应四个特征；当该状态出现时，这四个特征变为活动。
具体地说，特征向量 $\mathbf{x}(s)$ 在每个铺片中都有一个分量，因此共有 $4\times4\times4=64$ 个分量；除 $s$ 所落入的四个瓦片对应的分量之外，其余全为 $0$。
图 9.10 显示了在 1000 状态随机游走例子中，多个偏移铺片（粗编码）相对于单个铺片的优势。

瓦片编码的一个直接实用优势是，因为它使用划分，所以任一时刻活动特征的总数对任何状态都是相同的。
每个铺片中恰好有一个特征出现，因此出现特征的总数总等于铺片数。
这使步长参数 $\alpha$ 可以用简单直观的方式设置。
例如，若取 $\alpha=\frac{1}{n}$，其中 $n$ 是铺片数，就会得到精确的一次试验学习。
如果样例 $s\mapsto v$ 被训练，那么不管先前估计 $\hat{v}(s,\mathbf{w}_t)$ 是什么，新估计都会是 $\hat{v}(s,\mathbf{w}_{t+1})=v$。
通常我们希望变化比这慢一些，以便容纳泛化和目标输出中的随机变动。
例如，可以选择 $\alpha=\frac{1}{10n}$；此时被训练状态的估计会在一次更新中向目标移动十分之一，而邻近状态移动更少，移动量与它们共有的瓦片数成正比。

\begin{figure}[htbp]
\centering
\includegraphics[width=0.62\textwidth]{tex/强化学习/figures/math-10/figure-9-10.png}
\caption{为什么使用粗编码。图中给出了 1000 状态随机游走例子上，梯度蒙特卡洛算法使用单个铺片和多个铺片时的学习曲线。1000 个状态的空间被看作一个连续维度，用宽度为 200 个状态的瓦片覆盖。多个铺片彼此偏移 4 个状态。步长参数被设置成使两种情形的初始学习速率相同：单个铺片取 $\alpha=0.0001$，50 个铺片取 $\alpha=0.0001/50$。}
\end{figure}

瓦片编码还从二值特征向量的使用中获得计算优势。
因为每个分量不是 $0$ 就是 $1$，构成近似价值函数 (9.8) 的加权和几乎不需要计算。
不必执行 $d$ 次乘法和加法，只需计算 $n\ll d$ 个活动特征的索引，然后把权重向量中相应的 $n$ 个分量相加即可。

若其他状态落在任何相同瓦片中，就会从被训练状态向这些状态发生泛化，其强度与共有瓦片数成正比。
甚至如何让铺片彼此偏移的选择也会影响泛化。
如果铺片像图 9.9 中那样在每个维度上均匀偏移，那么不同状态可以以性质上不同的方式发生泛化，如图 9.11 上半部所示。
每个八个子图都显示了从一个被训练状态向附近点泛化的模式。
在这个例子中有八个铺片，因此一个瓦片内有 64 个子区域会以不同方式泛化，但都遵循这八种模式之一。
注意，均匀偏移会在许多模式中产生沿对角线的强效应。
如果铺片采用非对称偏移，如图下半部所示，就可以避免这些伪影。
这些下方的泛化模式更好，因为它们都很好地以被训练状态为中心，没有明显的不对称性。

\begin{figure}[htbp]
\centering
\includegraphics[width=0.9\textwidth]{tex/强化学习/figures/math-10/figure-9-11.png}
\caption{为什么瓦片编码中偏好非对称偏移。图中显示八个铺片时，从由小黑色加号标出的被训练状态到附近状态的泛化强度。若铺片均匀偏移（上方），则存在对角线伪影，并且泛化有显著变化；若铺片非对称偏移，则泛化更接近圆形且更均匀。}
\end{figure}

所有情形中的铺片都在每个维度上彼此偏移一个瓦片宽度的分数。
若 $w$ 表示瓦片宽度，$n$ 表示铺片数，那么 $\frac{w}{n}$ 是一个基本单位。
在边长为 $\frac{w}{n}$ 的小方块内，所有状态都激活相同瓦片，具有相同特征表示，因此有相同近似价值。
如果状态在任意笛卡尔方向上移动 $\frac{w}{n}$，特征表示就会改变一个分量或一个瓦片。
均匀偏移的铺片彼此正好偏移这个单位距离。
对二维空间，我们说每个铺片都按位移向量 $(1,1)$ 偏移，意思是它相对于前一铺片偏移 $\frac{w}{n}$ 倍的这个向量。
用这个术语说，图 9.11 下半部所示的非对称偏移铺片按位移向量 $(1,3)$ 偏移。

人们已经广泛研究过不同位移向量对瓦片编码泛化的影响（Parks and Militzer, 1991；An, 1991；An, Miller and Parks，
1991；Miller, An, Glanz and Carter, 1990），评估它们的均匀性以及产生像 $(1,1)$ 位移向量那样的对角线伪影的倾向。
基于这些工作，Miller 和 Glanz（1996）建议使用由最前面的奇数组成的位移向量。
特别地，对于维度为 $k$ 的连续空间，一个好的选择是使用前 $k$ 个奇数 $(1,3,5,7,\ldots,2k-1)$，并把 $n$（铺片数）设为大于 $4k$ 的 2 的整数次幂。
这就是我们在图 9.11 下半部生成铺片所采用的做法；其中 $k=2$，$n=2^3>4k$，位移向量为 $(1,3)$。
在三维情形中，前四个铺片相对于基准位置的总偏移将为 $(0,0,0)$、$(1,3,5)$、$(2,6,10)$ 和 $(3,9,15)$。
能够为任意 $k$ 高效生成这类铺片的开源软件已经很容易获得。

选择铺片策略时，必须选择铺片数以及瓦片形状。
铺片数与瓦片大小一起决定渐近近似的分辨率或精细程度，这与一般粗编码一样，也如图 9.8 所示。
瓦片形状会决定泛化的性质，如图 9.7 所示。
方形瓦片会在每个维度上大致均匀地泛化，如图 9.11 下半部所示。
沿一个维度拉长的瓦片，例如图 9.12 中图的条纹铺片，会促进沿该维度的泛化。
图 9.12 中图的铺片在左侧也更密、更细，因此会在该维度较小取值处促进水平方向上的区分。
图 9.12 右图中的对角条纹铺片会促进沿一个对角方向的泛化。
在更高维中，轴对齐条纹等价于在某些铺片中忽略若干维度，也就是超平面切片。
像图 9.12 左图这样的不规则铺片也是可能的，不过在标准软件之外的实践中较少见。

\begin{figure}[htbp]
\centering
\includegraphics[width=0.68\textwidth]{tex/强化学习/figures/math-10/figure-9-12.png}
\caption{铺片不必是网格。它们可以任意成形且非均匀，同时在许多情形中仍然计算高效。}
\end{figure}

在实践中，常常希望在不同铺片中使用不同形状的瓦片。
例如，可以使用一些竖直条纹铺片和一些水平条纹铺片。
这会鼓励沿任一维度的泛化。
不过，仅有条纹铺片时，不可能学到某个特定横纵坐标组合具有独特价值（无论为它学到什么，都会扩散到具有相同横坐标和纵坐标的状态）。
为此需要结合式的矩形瓦片，例如图 9.9 最初显示的那些。
使用多个铺片，其中一些水平、一些竖直、一些结合式，就能兼得两者：既偏好沿每个维度泛化，又能够学习组合处的特定价值（例子见 Sutton, 1996）。
铺片选择决定泛化；在这种选择能够被有效自动化之前，瓦片编码的重要之处在于它允许以灵活且符合人类直觉的方式作出选择。

另一个有助于降低内存需求的实用技巧是\textbf{哈希}，即把一个大的铺片以一致的伪随机方式折叠成一组小得多的瓦片。
哈希产生的瓦片由不相邻、不连通、随机散布在整个状态空间中的区域组成，但它们仍然形成一个穷尽性划分。
例如，一个瓦片可以由右侧所示的四个子瓦片组成。

\begin{figure}[htbp]
\centering
\includegraphics[width=0.24\textwidth]{tex/强化学习/figures/math-10/hash-tile-example.png}
\caption*{哈希后的一个瓦片可以由状态空间中分离的多个区域组成。}
\end{figure}

通过哈希，内存需求通常可以大幅降低，而性能损失很小。
这是可能的，因为只有状态空间的一小部分需要高分辨率。
哈希使我们在某种意义上摆脱维数灾难：内存需求不必随维度数指数增长，而只需匹配任务的实际需求。
瓦片编码的开源实现通常都包含高效哈希。

\noindent\textbf{练习 9.4}
假设我们相信两个状态维度中的一个比另一个更可能影响价值函数，并且泛化应当主要跨过这个维度，而不是沿着它发生。
可以使用什么样的铺片来利用这一先验知识？

\subsubsection*{9.5.5\quad 径向基函数}

\textbf{径向基函数}（radial basis functions, RBFs）是把粗编码自然推广到连续值特征的形式。
每个特征不再只能是 $0$ 或 $1$，而是可以取区间 $[0,1]$ 中的任意值，反映该特征出现的不同程度。
典型的 RBF 特征 $x_i$ 具有高斯（钟形）响应，$x_i(s)$ 依赖于状态 $s$ 与该特征的原型或中心状态 $c_i$ 之间的距离，以及该特征的相对宽度 $\sigma_i$：
\[
x_i(s)\doteq \exp\!\left(-\frac{\|s-c_i\|^2}{2\sigma_i^2}\right).
\]
当然，范数或距离度量可以按最适合当前状态和任务的方式选择。
下图显示了一个使用欧氏距离度量的一维例子。

\begin{figure}[htbp]
\centering
\includegraphics[width=0.55\textwidth]{tex/强化学习/figures/math-10/figure-9-13.png}
\caption{一维径向基函数。}
\end{figure}

相对于二值特征，RBF 的主要优点是它们产生的近似函数平滑变化且可微。
虽然这很有吸引力，但在多数情形中并没有实践上的重要性。
尽管如此，分级响应函数，例如 RBF，已经在瓦片编码背景下得到广泛研究（An, 1991；Miller et al., 1991；An et al., 1991；Lane, Handelman and Gelfand, 1992）。
所有这些方法都需要比瓦片编码多得多的计算复杂性，而且当状态维度超过两个时，常常会降低性能。
在高维中，瓦片边缘重要得多，而事实证明，很难在边缘附近得到控制良好的分级瓦片激活。

\textbf{RBF 网络}是使用 RBF 作为特征的线性函数近似器。
学习完全由公式 (9.7) 和 (9.8) 定义，就像其他线性函数近似器一样。
此外，有些 RBF 网络学习方法还会改变特征的中心和宽度，从而进入非线性函数近似器的范围。
非线性方法可能能够更精确地拟合目标函数。
RBF 网络，尤其是非线性 RBF 网络的缺点，是计算复杂性更高，并且在学习变得稳健高效之前，往往需要更多手工调参。

\subsection*{9.6\quad 手动选择步长参数}

大多数 SGD 方法都要求设计者选择合适的步长参数 $\alpha$。
理想情况下，这种选择应当自动化；在某些情形中确实如此，但在多数情形中，手动设置仍然是常见做法。
为了做到这一点，并更好地理解算法，形成一些关于步长参数作用的直觉是有用的。
一般而言，我们能否说明它应当如何设置？

遗憾的是，理论考虑帮助不大。
随机近似理论给出了关于缓慢递减步长序列的条件 (2.7)，这些条件足以保证收敛，但往往导致学习过慢。
经典选择 $\alpha_t=1/t$ 会在表格型 MC 方法中产生样本平均，但它不适合 TD 方法、非平稳问题，或任何使用函数近似的方法。
对于线性方法，有一些递归最小二乘方法可以设置最优\textbf{矩阵}步长；这些方法也可以扩展到时序差分学习，例如第 9.8 节描述的 LSTD 方法。
但它们需要 $O(d^2)$ 个步长参数，也就是比我们正在学习的参数多 $d$ 倍。
因此，在最需要函数近似的大规模问题中，我们排除使用它们。

为了获得关于如何手动设置步长参数的直觉，最好暂时回到表格型情形。
在那里我们可以理解，步长 $\alpha=1$ 会使一个目标之后的样本误差被完全消除（见 (2.4)，步长取 1）。
如第 201 页所讨论的，我们通常希望学习比这更慢。
在表格型情形中，步长 $\alpha=\frac{1}{10}$ 会让约 10 次经验后大致收敛到它们的平均目标；如果我们希望用 100 次经验学习，就用 $\alpha=\frac{1}{100}$。
一般来说，如果 $\alpha=\frac{1}{\tau}$，那么一个状态的表格估计将在大约 $\tau$ 次该状态经验后接近其目标的均值，且最近目标影响最大。

在一般函数近似中，状态经验的\textbf{次数}没有这么清楚的概念，因为每个状态都可能以不同程度与所有其他状态相似和不相似。
不过，在线性函数近似情形中有一个类似规则，会给出类似行为。
假设你希望用大约 $\tau$ 次具有实质上相同特征向量的经验完成学习。
设置线性 SGD 方法步长参数的一个好经验规则是
\begin{equation}
\alpha\doteq \left(\tau\mathbb{E}[\mathbf{x}^\top\mathbf{x}]\right)^{-1},
\tag{9.19}
\end{equation}
其中 $\mathbf{x}$ 是从与 SGD 中输入向量相同分布中选择的随机特征向量。
如果特征向量长度变化不大，这个方法效果最好；理想情况下 $\mathbf{x}^\top\mathbf{x}$ 是常数。

\noindent\textbf{练习 9.5}
假设你使用瓦片编码把七维连续状态空间转换为二值特征向量，以估计状态价值函数 $\hat{v}(s,\mathbf{w})\approx v_\pi(s)$。
你认为各维度之间没有强交互，因此决定对每个维度分别使用八个铺片（条纹铺片），共 $7\times8=56$ 个铺片。
此外，为了防备维度之间存在一些两两交互，你又取所有 $\binom{7}{2}=21$ 对维度，并对每一对维度用矩形瓦片作结合式铺片。
你为每对维度做两个铺片，总计 $21\times2+56=98$ 个铺片。
给定这些特征向量，你怀疑仍需平均掉一些噪声，因此决定让学习较为渐进，希望在具有相同特征向量的约 10 次出现之后，学习接近其渐近值。
应当使用什么步长参数 $\alpha$？
为什么？

\subsection*{9.7\quad 非线性函数近似：人工神经网络}

人工神经网络（artificial neural networks, ANNs）被广泛用于非线性函数近似。
ANN 是由互连单元组成的网络，这些单元具有神经元的一些性质；神经元是神经系统的主要组成部分。
ANN 有很长的历史，而训练深层 ANN（深度学习）的最新进展，已经使机器学习系统，包括强化学习系统，具备了一些最令人印象深刻的能力。
第 16 章将描述若干使用 ANN 函数近似的强化学习系统的精彩例子。

图 9.14 显示了一个一般的前馈 ANN；前馈意味着网络中没有环，也就是说，网络中不存在某个单元的输出能够影响其输入的路径。
图中的网络有一个由两个输出单元组成的输出层、一个由四个输入单元组成的输入层，以及两个“隐藏层”，即既不是输入层也不是输出层的层。
每条连接都关联一个实值权重。
权重大致对应于真实神经网络中突触连接的效能（见第 15.1 节）。
如果一个 ANN 的连接中至少有一个环，那么它就是循环 ANN，而不是前馈 ANN。
虽然前馈 ANN 和循环 ANN 都曾用于强化学习，这里我们只考察更简单的前馈情形。

\begin{figure}[htbp]
\centering
\includegraphics[width=0.55\textwidth]{tex/强化学习/figures/math-10/figure-9-14.png}
\caption{一个通用前馈 ANN，具有四个输入单元、两个输出单元和两个隐藏层。}
\end{figure}

单元（图 9.14 中的圆）通常是半线性单元，也就是说，它们先计算输入信号的加权和，然后对结果应用非线性函数，称为\textbf{激活函数}，以产生该单元的输出或激活。
会使用不同的激活函数，但它们通常是 S 形或 sigmoid 函数，例如 logistic 函数 $f(x)=1/(1+e^{-x})$；不过有时也使用 rectifier 非线性 $f(x)=\max(0,x)$。
像 $x\ge\theta$ 时 $f(x)=1$、否则为 $0$ 的阶跃函数，会产生阈值为 $\theta$ 的二值单元。
网络输入层中的单元略有不同，因为它们的激活被设为外部给定的数值，即该网络正在近似的函数的输入。

前馈 ANN 每个输出单元的激活，都是网络输入单元激活模式的非线性函数。
这些函数由网络的连接权重参数化。
没有隐藏层的 ANN 只能表示所有可能输入--输出函数中的极小一部分。
然而，具有单个隐藏层且包含足够多有限个 sigmoid 单元的 ANN，可以把网络输入空间某个紧集上的任意连续函数近似到任意精度（Cybenko, 1989）。
对满足温和条件的其他非线性激活函数也是如此，但非线性是必要的：如果多层前馈 ANN 中所有单元都具有线性激活函数，那么该网络等价于没有隐藏层的网络（因为线性函数的线性函数仍然是线性的）。

尽管单隐藏层 ANN 具有这种“通用近似”性质，经验和理论都表明，对许多人工智能任务所需的复杂函数进行近似，如果有由许多层较低层抽象层级组合而成的抽象，会容易得多，甚至可能必须如此；也就是说，需要由深层架构（例如有许多隐藏层的 ANN）产生的抽象。（详尽综述见 Bengio, 2009。）
深层 ANN 的相继各层会计算越来越
抽象的网络“原始”输入表示，每个单元都提供一个特征，共同贡献于网络整体输入--输出函数的层次化表示。

因此，训练 ANN 的隐藏层，是一种为给定问题自动创造合适特征的方式，使层次化表示可以在不完全依赖手工设计特征的情况下产生。
这一直是人工智能面临的持久挑战，也解释了为什么多年来带隐藏层 ANN 的学习算法受到如此多关注。
ANN 通常通过随机梯度方法学习（第 9.3 节）。
每个权重都沿着旨在改善网络整体性能的方向调整；整体性能由某个需要最小化或最大化的目标函数度量。
在最常见的监督学习情形中，目标函数是一组带标签训练样例上的期望误差或损失。
在强化学习中，ANN 可以使用 TD 误差来学习价值函数，也可以像梯度赌博机（第 2.8 节）或策略梯度算法（第 13 章）那样，旨在最大化期望奖励。
在所有这些情形中，都必须估计每个连接权重的变化会如何影响网络整体性能；换句话说，在给定网络所有权重当前值的情况下，估计目标函数关于每个权重的偏导数。
梯度就是这些偏导数组成的向量。

对带隐藏层 ANN 来说，最成功的做法（前提是单元的激活函数可微）是反向传播算法；它由在网络中交替进行的前向和反向传播组成。
每次前向传播在给定网络输入单元当前激活的情况下，计算每个单元的激活。
每次前向传播之后，一次反向传播会高效地为每个权重计算一个偏导数。
（和其他随机梯度学习算法一样，这些偏导数组成的向量是真实梯度的一个估计。）
第 15.10 节将讨论一些用强化学习原理而非反向传播来训练带隐藏层 ANN 的方法。
这些方法不如反向传播算法高效，但可能更接近真实神经网络的学习方式。

反向传播算法能够在具有 1 个或 2 个隐藏层的浅层网络上产生良好结果，但对更深的 ANN 可能效果不好。
事实上，训练一个有 $k+1$ 个隐藏层的网络，可能比训练一个有 $k$ 个隐藏层的网络表现更差，尽管更深网络能够表示浅层网络能够表示的所有函数（Bengio, 2009）。
解释这类结果并不容易，但有几个因素很重要。
首先，典型深层 ANN 中大量权重使过拟合问题难以避免；过拟合是指网络无法正确泛化到未训练过的情形。
其次，反向传播不适合深层 ANN，因为它的反向传播所计算出的偏导数要么会在向输入端传播时迅速衰减，使深层学习极其缓慢；要么会在向输入端传播时迅速增大，使学习不稳定。
处理这些问题的方法，在很大程度上促成了使用深层 ANN 的系统近年来取得的许多令人印象深刻的结果。

对于任何基于有限训练数据、调整具有许多自由度函数的函数近似方法，过拟合都是问题。
对于不依赖有限训练集的在线强化学习来说，它的问题较小，但有效泛化仍然是重要问题。
过拟合是 ANN 的一般问题，对深层 ANN 尤其如此，因为它们往往有非常多的权重。
已经发展出许多减少过拟合的方法。
这些方法包括：当验证数据（交叉验证）上的性能开始下降时停止训练；修改目标函数以惩罚近似的复杂性（正则化）；以及引入权重之间的依赖关系以减少自由度数量（例如权重共享）。

减少深层 ANN 过拟合的一种特别有效方法，是 Srivastava、Hinton、Krizhevsky、Sutskever 和 Salakhutdinov（2014）提出的 dropout 方法。
训练期间，网络中的单元会连同其连接一起被随机移除（drop out）。
这可以看作训练大量“变薄”的网络。
在测试时组合这些变薄网络的结果，是一种改善泛化性能的方法。
dropout 方法通过把每个单元的出边权重乘以该单元在训练中被保留的概率，高效地近似这种组合。
Srivastava 等人发现，该方法显著改善了泛化性能。
它鼓励单个隐藏单元学习能够与其他特征的随机集合良好配合的特征。
这增加了隐藏单元形成的特征的通用性，使网络不会过度专门化于罕见情形。

Hinton、Osindero 和 Teh（2006）在解决深层 ANN 深层训练问题方面迈出了重要一步；他们使用深度置信网络，这是一类与这里讨论的深层 ANN 密切相关的网络。
在他们的方法中，最深层使用无监督学习算法一次训练一层。
无监督学习不依赖整体目标函数，可以提取捕捉输入流统计规律的特征。
最深层先被训练；然后用这个已训练层提供输入，训练下一深层，依此类推，直到网络所有或许多层的权重被设置为可作为监督学习初始值的数值。
随后，网络再通过反向传播相对于整体目标函数进行微调。
研究表明，这种方法通常明显优于从随机值初始化权重后进行反向传播。
以这种方式初始化权重的网络表现更好可能有许多原因；一种想法是，该方法把网络置于权重空间的某个区域，使基于梯度的算法能够取得良好进展。

\textbf{批量归一化}（Ioffe and Szegedy, 2015）是另一种使深层 ANN 更容易训练的技术。
人们早就知道，如果网络输入被归一化，例如把每个输入变量调整为零均值和单位方差，那么 ANN 学习会更容易。
用于训练深层 ANN 的批量归一化，会在深层的输出进入下一层之前对其进行归一化。
Ioffe 和 Szegedy（2015）使用训练样例的子集，即“小批量”的统计量，来归一化这些层间信号，从而改善深层 ANN 的学习速率。

训练深层 ANN 的另一种有用技术是\textbf{深度残差学习}（He, Zhang, Ren, and Sun, 2016）。
有时，学习一个函数与恒等函数之间的差异，比学习该函数本身更容易。
然后把这个差异，即残差函数，加到输入上，就得到所需函数。
在深层 ANN 中，可以让一组层通过在该组周围添加快捷连接或跳跃连接来学习残差函数。
这些连接把该组的输入加到输出上，并且不需要额外权重。
He 等人（2016）用在每一对相邻层周围都带跳跃连接的深度卷积网络评估了这种方法，发现它在基准图像分类任务上相比没有跳跃连接的网络有显著改进。
批量归一化和深度残差学习都用于第 16 章将描述的围棋强化学习应用。

一种在包括令人印象深刻的强化学习应用（第 16 章）在内的应用中被证明非常成功的深层 ANN，是\textbf{深度卷积网络}。
这类网络专门处理以空间阵列排列的高维数据，例如图像。
它受到大脑早期视觉处理方式的启发（LeCun, Bottou, Bengio and Haffner, 1998）。
由于其特殊结构，深度卷积网络可以通过反向传播训练，而无需求助于上述训练深层的方法。

图 9.15 展示了一个深度卷积网络的结构。
这个来自 LeCun 等人（1998）的例子被设计用于识别手写字符。
它由交替的卷积层和子采样层组成，随后是若干全连接最终层。
每个卷积层产生若干特征图。
特征图是在一组单元上的活动模式；其中每个单元都对其感受野中的数据执行相同操作，感受野是它从前一层（或者在第一卷积层中从外部输入）“看到”的数据部分。
一个特征图中的单元除了感受野不同外彼此相同；这些感受野大小和形状都相同，只是平移到输入数据阵列的不同位置。
同一特征图中的单元共享相同权重。
这意味着无论特征位于输入阵列的哪个位置，特征图都会检测同一个特征。
例如，在图 9.15 的网络中，第一卷积层产生 6 个特征图，每个特征图由 $28\times28$ 个单元组成。
每个特征图中的每个单元都有 $5\times5$ 的感受野，并且这些感受野彼此重叠（本例中重叠四列和四行）。
因此，6 个特征图中的每一个都只由 25 个可调权重指定。

\begin{figure}[htbp]
\centering
\includegraphics[width=0.9\textwidth]{tex/强化学习/figures/math-10/figure-9-15.png}
\caption{深度卷积网络。经 IEEE Proceedings 授权重印，来自 LeCun、Bottou、Bengio 和 Haffner 的“Gradient-based learning applied to document recognition”，第 86 卷，1998；授权通过 Copyright Clearance Center, Inc. 转达。}
\end{figure}

深度卷积网络的子采样层会降低特征图的空间分辨率。
子采样层中的每个特征图由若干单元组成，这些单元会在前一卷积层特征图中的一组单元感受野上取平均。
例如，图 9.15 网络第一个子采样层中，6 个特征图中的每个单元都会在卷积层产生的一个特征图上，对一个 $2\times2$ 非重叠感受野取平均，从而得到 14 个 $14\times14$ 的特征图。
子采样层降低了网络对所检测特征空间位置的敏感性，也就是说，它们帮助使网络响应具有空间不变性。
这很有用，因为在图像中某处检测到的特征，也很可能在其他位置有用。

ANN 设计和训练方面的进展（这里只提到其中少数）都对强化学习有贡献。
虽然当前强化学习理论主要局限于使用表格型或线性函数近似的方法，但著名强化学习应用的出色表现，很大程度上归功于多层 ANN 的非线性函数近似。
第 16 章将讨论其中几个应用。

\subsection*{9.8\quad 最小二乘 TD}

本章到目前为止讨论的所有方法，每个时间步的计算量都与参数个数成正比。
不过，如果使用更多计算，就可以做得更好。
本节介绍一种用于线性函数近似的方法，可以说是在这种情形下能做到的最好方法。

正如第 9.4 节所确立的，带线性函数近似的 TD(0) 会在渐近意义下（若步长适当地递减）收敛到 TD 不动点：
\[
\mathbf{w}_{\mathrm{TD}}=\mathbf{A}^{-1}\mathbf{b},
\]
其中
\[
\mathbf{A}\doteq\mathbb{E}\!\left[\mathbf{x}_t(\mathbf{x}_t-\gamma\mathbf{x}_{t+1})^\top\right],
\qquad
\mathbf{b}\doteq\mathbb{E}[R_{t+1}\mathbf{x}_t].
\]
有人可能会问，为什么必须迭代地计算这个解？
这太浪费数据了。
难道不能通过计算 $\mathbf{A}$ 和 $\mathbf{b}$ 的估计，然后直接计算 TD 不动点来做得更好吗？
\textbf{最小二乘 TD} 算法，通常称为 LSTD，正是这样做的。
它形成自然估计
\begin{equation}
\hat{\mathbf{A}}_t\doteq
\sum_{k=0}^{t-1}\mathbf{x}_k(\mathbf{x}_k-\gamma\mathbf{x}_{k+1})^\top+\varepsilon\mathbf{I},
\qquad
\hat{\mathbf{b}}_t\doteq
\sum_{k=0}^{t-1}R_{k+1}\mathbf{x}_k,
\tag{9.20}
\end{equation}
其中 $\mathbf{I}$ 是单位矩阵，而对某个小的 $\varepsilon>0$，$\varepsilon\mathbf{I}$ 保证 $\hat{\mathbf{A}}_t$ 始终可逆。
看起来这些估计似乎都应除以 $t$，事实上也确实应当如此；但按这里的定义，它们实际上是 $t$ 倍的 $\mathbf{A}$ 和 $t$ 倍的 $\mathbf{b}$ 的估计。
不过，当 LSTD 用这些估计来估计 TD 不动点时，额外的 $t$ 因子相互抵消：
\begin{equation}
\mathbf{w}_t\doteq \hat{\mathbf{A}}_t^{-1}\hat{\mathbf{b}}_t.
\tag{9.21}
\end{equation}

这个算法是线性 TD(0) 数据效率最高的形式，但计算代价也更高。
回忆一下，半梯度 TD(0) 只需要 $O(d)$ 的内存和每步计算量。

LSTD 有多复杂？
按照上面写法，其复杂度似乎会随 $t$ 增长，但 (9.20) 中两个近似量可以用我们之前介绍的技术（例如第 2 章中的技术）增量实现，使每一步都能在常数时间完成。
即便如此，$\hat{\mathbf{A}}_t$ 的更新涉及一个外积（列向量乘行向量），因此是一次矩阵更新；其计算复杂度为 $O(d^2)$，当然保存 $\hat{\mathbf{A}}_t$ 矩阵所需内存也是 $O(d^2)$。

一个潜在更大的问题是，最终计算 (9.21) 要用到 $\hat{\mathbf{A}}_t$ 的逆，而一般求逆的计算复杂度是 $O(d^3)$。
幸运的是，我们这种特殊形式的矩阵（外积之和）的逆也可以只用 $O(d^2)$ 次计算来增量更新：
\begin{align}
\hat{\mathbf{A}}_t^{-1}
&=\left(\hat{\mathbf{A}}_{t-1}+\mathbf{x}_t(\mathbf{x}_t-\gamma\mathbf{x}_{t+1})^\top\right)^{-1}
\qquad\text{（由 (9.20)）}\notag\\
&=\hat{\mathbf{A}}_{t-1}^{-1}
-\frac{\hat{\mathbf{A}}_{t-1}^{-1}\mathbf{x}_t(\mathbf{x}_t-\gamma\mathbf{x}_{t+1})^\top\hat{\mathbf{A}}_{t-1}^{-1}}
{1+(\mathbf{x}_t-\gamma\mathbf{x}_{t+1})^\top\hat{\mathbf{A}}_{t-1}^{-1}\mathbf{x}_t},
\tag{9.22}
\end{align}
其中 $t>0$，且 $\hat{\mathbf{A}}_0\doteq \varepsilon\mathbf{I}$。
虽然恒等式 (9.22)，即 Sherman--Morrison 公式，表面上复杂，但它只涉及向量--矩阵和向量--向量乘法，因此只有 $O(d^2)$。
于是，我们可以存储逆矩阵 $\hat{\mathbf{A}}_t^{-1}$，用 (9.22) 维护它，然后在 (9.21) 中使用它，所有这些都只需 $O(d^2)$ 内存和每步计算量。
完整算法见下一页的框。

当然，$O(d^2)$ 仍然显著高于半梯度 TD 的 $O(d)$。
LSTD 更高的数据效率是否值得这种计算代价，取决于 $d$ 有多大、快速学习有多重要，以及系统其他部分的代价。
有时人们还会指出 LSTD 不需要步长参数，但这个优点可能被夸大了。
LSTD 不需要步长，但它需要 $\varepsilon$；如果 $\varepsilon$ 选得太小，逆矩阵序列可能剧烈变化；如果 $\varepsilon$ 选得太大，则学习会变慢。
此外，LSTD 缺少步长参数意味着它永远不会遗忘。
这有时是可取的，但如果目标策略 $\pi$ 像强化学习和 GPI 中那样变化，就会有问题。
在控制应用中，LSTD 通常必须与其他机制结合以引入遗忘，从而抵消不需要步长参数的任何初始优势。

\begin{center}
\fbox{\begin{minipage}{0.9\textwidth}
\noindent\textbf{LSTD，用于估计 $\hat{v}=\mathbf{w}^\top\mathbf{x}(\cdot)\approx v_\pi$（$O(d^2)$ 版本）}

\medskip
\noindent 输入：特征表示 $\mathbf{x}:\mathcal{S}^+\to\mathbb{R}^d$，且 $\mathbf{x}(\mathrm{terminal})=\mathbf{0}$。

\noindent 算法参数：小的 $\varepsilon>0$。

\medskip
\noindent $\hat{\mathbf{A}}^{-1}\leftarrow \varepsilon^{-1}\mathbf{I}$\hfill 一个 $d\times d$ 矩阵

\noindent $\hat{\mathbf{b}}\leftarrow \mathbf{0}$\hfill 一个 $d$ 维向量

\noindent 对每一幕循环：

\quad 初始化 $S$；$\mathbf{x}\leftarrow \mathbf{x}(S)$。

\quad 对幕中的每一步循环：

\qquad 选择并采取动作 $A\sim\pi(\cdot|S)$，观察 $R,S'$；$\mathbf{x}'\leftarrow \mathbf{x}(S')$。

\qquad $\mathbf{v}\leftarrow \hat{\mathbf{A}}^{-1\top}(\mathbf{x}-\gamma\mathbf{x}')$。

\qquad $\hat{\mathbf{A}}^{-1}\leftarrow
\hat{\mathbf{A}}^{-1}-(\hat{\mathbf{A}}^{-1}\mathbf{x})\mathbf{v}^\top/(1+\mathbf{v}^\top\mathbf{x})$。

\qquad $\hat{\mathbf{b}}\leftarrow \hat{\mathbf{b}}+R\mathbf{x}$。

\qquad $\mathbf{w}\leftarrow \hat{\mathbf{A}}^{-1}\hat{\mathbf{b}}$。

\qquad $S\leftarrow S'$；$\mathbf{x}\leftarrow \mathbf{x}'$。

\quad 直到 $S'$ 为终止状态。
\end{minipage}}
\end{center}

\subsection*{9.9\quad 基于记忆的函数近似}

到目前为止，我们讨论的是近似价值函数的\textbf{参数化}方法。
在这种方法中，学习算法调整某个函数形式的参数，而这个函数形式旨在近似整个问题状态空间上的价值函数。
每次更新 $s\mapsto g$ 都是一个训练样例，学习算法使用它来改变参数，以减少近似误差。
更新之后，训练样例可以被丢弃（虽然也可以保存下来以后再次使用）。
当需要某个状态（我们称为\textbf{查询状态}）的近似价值时，只需用学习算法产生的最新参数在该状态处评价函数。

基于记忆的函数近似方法非常不同。
它们只是把训练样例到达时保存到记忆中（或者至少保存样例的一个子集），而不更新任何参数。
然后，只要需要某个查询状态的价值估计，就从记忆中取出一组样例，并用它们计算该查询状态的价值估计。
这种方法有时称为\textbf{惰性学习}，因为对训练样例的处理被推迟到系统被查询、需要给出输出时才进行。

基于记忆的函数近似方法是\textbf{非参数}方法的典型例子。
不同于参数化方法，近似函数的形式不被限制在固定的参数化函数类中，例如线性函数或多项式；它由训练样例本身以及某种组合它们以输出查询状态估计价值的方式决定。
随着越来越多训练样例在记忆中积累，人们期望非参数方法能对任意目标函数产生越来越准确的近似。

基于记忆的方法有许多不同形式，取决于存储的训练样例如何选择，以及它们如何用于响应查询。
这里我们关注\textbf{局部学习}方法，也就是只在当前查询状态附近局部近似价值函数的方法。
这些方法从记忆中取回一组训练样例；这些样例的状态被认为与查询状态最相关。
相关性通常取决于状态之间的距离：一个训练样例的状态越接近查询状态，就越相关；距离可以用许多不同方式定义。
在给出查询状态的价值后，局部近似被丢弃。

基于记忆方法最简单的例子是\textbf{最近邻}方法。
它只需在记忆中找到状态最接近查询状态的那个样例，并把该样例的价值作为查询状态的近似价值返回。
换句话说，如果查询状态为 $s$，而 $s'\mapsto g$ 是记忆中状态 $s'$ 最接近 $s$ 的样例，那么返回 $g$ 作为 $s$ 的近似价值。
稍复杂一些的是\textbf{加权平均}方法；它取回一组最近邻样例，并返回其目标值的加权平均，其中权重通常随样例状态与查询状态之间距离增加而减小。
\textbf{局部加权回归}类似，但它用一种参数化近似方法，在一组最近状态的价值上拟合一个曲面；该方法最小化类似 (9.1) 的加权误差度量，权重取决于到查询状态的距离。
返回的值是在查询状态处评价这个局部拟合曲面，随后该局部近似曲面被丢弃。

由于是非参数的，基于记忆的方法相对于参数化方法的优点，是不把近似限制在预先指定的函数形式上。
这使精度可以随数据积累而提高。
基于记忆的局部近似方法还有其他性质，使它们很适合强化学习。
由于轨迹采样在强化学习中非常重要（见第 8.6 节），基于记忆的局部方法可以把函数近似集中在真实或模拟轨迹访问过的状态（或状态--动作对）的局部邻域。
可能不需要全局近似，因为状态空间中的许多区域永远（或几乎永远）不会到达。
此外，与参数化方法需要逐步调整全局近似的参数不同，基于记忆的方法允许智能体经验对当前状态邻域中的价值估计产生相对即时的影响。

避免全局近似也是应对维数灾难的一种方式。
例如，对具有 $k$ 个维度的状态空间，存储全局近似的表格型方法需要关于 $k$ 指数级的内存。
另一方面，在基于记忆的方法中，每个样例所需内存与 $k$ 成正比；若存储 $n$ 个样例，则所需内存关于 $n$ 线性。
没有任何东西关于 $k$ 或 $n$ 指数增长。
当然，关键的剩余问题是，基于记忆的方法能否足够快地回答查询，从而对智能体有用。
相关担忧是随着记忆规模增长，速度会如何下降。
在许多应用中，在大数据库中寻找最近邻可能太慢而不实用。

基于记忆方法的支持者发展了多种加速最近邻搜索的方法。
使用并行计算机或专用硬件是一种办法；另一种是使用专门的多维数据结构来存储训练数据。
为此应用研究过的一种数据结构是 $k$-d 树（即 $k$ 维树的缩写），它递归地把 $k$ 维空间划分为若干区域，并把这些区域组织成二叉树的节点。
根据数据量及其在状态空间中的分布，使用 $k$-d 树的最近邻搜索可以在搜索邻居时快速排除大区域，从而使一些朴素搜索过慢的问题变得可行。

局部加权回归还需要快速执行局部回归计算，而这种计算必须为每次查询重复。
研究人员已经发展出许多方法来处理这些问题，包括遗忘条目的方法，以把数据库大小保持在界内。
本章末尾的书目与历史评论部分会指向一些相关文献，其中包括若干描述基于记忆学习在强化学习中应用的论文。

\subsection*{9.10\quad 基于核的函数近似}

基于记忆的方法，例如上面描述的加权平均和局部加权回归，依赖于给数据库中的样例 $s'\mapsto g$ 分配权重，而权重取决于 $s'$ 与查询状态 $s$ 之间的距离。
分配这些权重的函数称为\textbf{核函数}，或简称\textbf{核}。
在上面的加权平均和局部加权回归方法中，例如，核函数 $k:\mathbb{R}\to\mathbb{R}$ 给状态间距离分配权重。
更一般地，权重不必依赖于距离；它们可以依赖于状态之间的某种其他相似性度量。
在这种情况下，$k:\mathcal{S}\times\mathcal{S}\to\mathbb{R}$，因此 $k(s,s')$ 是在回答关于 $s$ 的查询时，赋予关于 $s'$ 的数据的权重。

稍换一个角度看，$k(s,s')$ 是从 $s'$ 到 $s$ 的泛化强度的度量。
核函数以数值方式表达关于任意状态的知识与任意其他状态有多相关。
例如，图 9.11 中瓦片编码的泛化强度，对应于由均匀和非对称瓦片偏移产生的不同核函数。
虽然瓦片编码在操作中并不显式使用核函数，但它确实按照某个核来泛化。
事实上，正如下面将进一步讨论的，线性参数化函数近似所产生的泛化强度总能由核函数描述。

\textbf{核回归}是一种基于记忆的方法，它对记忆中所有样例的目标作核加权平均，并把结果赋给查询状态。
若 $\mathcal{D}$ 是存储样例集合，$g(s')$ 表示某个存储样例中状态 $s'$ 的目标，那么核回归根据 $\mathcal{D}$ 近似目标函数，在这里即价值函数：
\begin{equation}
\hat{v}(s,\mathcal{D})=\sum_{s'\in\mathcal{D}}k(s,s')g(s').
\tag{9.23}
\end{equation}
上面描述的加权平均方法是一个特殊情形，其中只有当 $s$ 和 $s'$ 足够接近时，$k(s,s')$ 才非零，因此求和不必遍历整个 $\mathcal{D}$。

常见核之一是第 9.5.5 节 RBF 函数近似中使用的高斯径向基函数（RBF）。
在那里描述的方法中，RBF 是特征，其中心和宽度要么从一开始固定（中心大概集中在预计会出现许多样例的区域），要么在学习中以某种方式调整。
除了会调整中心和宽度的方法外，这是一种线性参数化方法；其参数是每个 RBF 的权重，通常由随机梯度或半梯度下降学习。
近似形式是预先确定的 RBF 的线性组合。
使用 RBF 核的核回归与此有两点不同。
第一，它是基于记忆的：RBF 以存储样例的状态为中心。
第二，它是非参数的：没有要学习的参数；查询的响应由 (9.23) 给出。

当然，核回归的实际实现还必须处理许多问题，这些问题超出了我们简短讨论的范围。
不过事实证明，任何类似第 9.4 节所描述的线性参数化回归方法，只要状态由特征向量
\[
\mathbf{x}(s)=(x_1(s),x_2(s),\ldots,x_d(s))^\top
\]
表示，都可以重新表述为核回归，其中 $k(s,s')$ 是 $s$ 与 $s'$ 的特征向量表示的内积，即
\begin{equation}
k(s,s')=\mathbf{x}(s)^\top\mathbf{x}(s').
\tag{9.24}
\end{equation}
使用这个核函数的核回归，会产生与线性参数化方法在使用这些特征向量并用相同训练数据学习时相同的近似。

我们略去对此的数学论证；它可以在任何现代机器学习教材中找到，例如 Bishop（2006）。
这里只指出一个重要含义。
与其为线性参数化函数近似器构造特征，不如直接构造核函数，而完全不引用特征向量。
并非所有核函数都能像 (9.24) 那样表示为特征向量内积，但可以这样表示的核函数，能够相对于等价的参数化方法提供显著优势。
对许多特征向量集合，(9.24) 有紧凑的函数形式，可以在不进行 $d$ 维特征空间中任何计算的情况下求值。
在这些情况下，核回归比直接使用以这些特征向量表示状态的线性参数化方法复杂度低得多。
这就是所谓“核技巧”：它允许我们在一个扩展特征空间的高维中有效工作，而实际只处理存储训练样例的集合。
核技巧是许多机器学习方法的基础；研究人员已经展示了它有时如何有益于强化学习。

\subsection*{9.11\quad 更深入地看同策略学习：兴趣与强调}

本章到目前为止考察的算法，都把遇到的所有状态等同对待，仿佛它们都同等重要。
然而，在某些情形中，我们对一些状态比对另一些状态更感兴趣。
例如，在折扣分幕问题中，与后期状态相比，我们可能更关心准确评价一幕中的早期状态，因为折扣可能已经使后期奖励对起始状态价值的重要性小得多。
或者，如果正在学习动作价值函数，那么准确评价很差动作可能不那么重要，因为这些动作的价值远低于贪婪动作。
函数近似资源总是有限的，如果更有针对性地使用它们，就可能改善性能。

我们把遇到的所有状态等同对待的一个原因是：此时我们按照同策略分布进行更新；对于半梯度方法，这有很强的理论结果。
回忆一下，同策略分布定义为在遵循目标策略时 MDP 中遇到的状态分布。
现在我们将显著推广这个概念。
不再为 MDP 只有一个同策略分布，而是会有许多个。
它们的共同点是，都是遵循目标策略时轨迹中遇到的状态分布；不同点在于这些轨迹在某种意义上如何被初始化。

现在介绍一些新概念。
首先，引入一个非负标量度量，即称为\textbf{兴趣}的随机变量 $I_t$，表示我们对准确评价时刻 $t$ 的状态（或状态--动作对）有多感兴趣。
如果我们完全不关心这个状态，那么兴趣应为零；如果完全关心，它可以为一，不过形式上允许取任意非负值。
兴趣可以按任何因果方式设置；例如，它可以依赖于截至时刻 $t$ 的轨迹，或时刻 $t$ 的已学参数。
于是，$\overline{\mathrm{VE}}$ (9.1) 中的分布 $\mu$ 被定义为遵循目标策略时遇到的状态分布，并按兴趣加权。
其次，我们引入另一个非负标量随机变量，即\textbf{强调} $M_t$。
这个标量乘以学习更新，从而强调或弱化时刻 $t$ 所做的学习。
替代 (9.15) 的一般 $n$ 步学习规则为
\begin{equation}
\mathbf{w}_{t+n}\doteq
\mathbf{w}_{t+n-1}
+\alpha M_t\left[G_{t:t+n}-\hat{v}(S_t,\mathbf{w}_{t+n-1})\right]\nabla\hat{v}(S_t,\mathbf{w}_{t+n-1}),
\quad 0\le t<T,
\tag{9.25}
\end{equation}
其中 $n$ 步回报由 (9.16) 给出，强调由兴趣递归决定：
\begin{equation}
M_t=I_t+\gamma^n M_{t-n},
\quad 0\le t<T,
\tag{9.26}
\end{equation}
并且对所有 $t<0$，$M_t\doteq 0$。
这些方程也包括蒙特卡洛情形；此时 $G_{t:t+n}=G_t$，所有更新都在一幕结束时进行，$n=T-t$，且 $M_t=I_t$。
例 9.4 说明兴趣和强调如何产生更准确的价值估计。

\begin{center}
\fbox{\begin{minipage}{0.9\textwidth}
\noindent\textbf{例 9.4：兴趣与强调}

\medskip
为了看到使用兴趣和强调的潜在好处，考虑如下所示的四状态马尔可夫奖励过程：

\begin{center}
\includegraphics[width=0.78\textwidth]{tex/强化学习/figures/math-10/example-9-4-mrp.png}

\small\textbf{图示：}四状态马尔可夫奖励过程。
\end{center}

每一幕从最左状态开始，然后每一步以奖励 $+1$ 向右转移一个状态，直到到达终止状态。
因此，第一个状态的真实价值为 4，第二个状态为 3，依此类推，如每个状态下方所示。
这些是真实价值；由于参数化约束，估计价值只能近似它们。
参数向量有两个分量 $\mathbf{w}=(w_1,w_2)^\top$，参数化写在每个状态内部。
前两个状态的估计价值都只由 $w_1$ 给出，因此必须相同，尽管它们的真实价值不同。
类似地，第三和第四状态的估计价值都只由 $w_2$ 给出，也必须相同，尽管它们的真实价值不同。
假设我们只对准确评价最左状态感兴趣；我们为它分配兴趣 1，而所有其他状态的兴趣为 0，如状态上方所示。

先考虑把梯度蒙特卡洛算法应用于这个问题。
本章前面给出的、不考虑兴趣和强调的方法（(9.7) 以及第 202 页框中的算法）会在步长递减时收敛到参数向量 $\mathbf{w}_\infty=(3.5,1.5)$；这给第一个状态，也就是唯一让我们感兴趣的状态，赋值 3.5（即第一和第二状态真实价值的中间值）。
另一方面，本节给出的使用兴趣和强调的方法，会精确正确地学到第一个状态的价值；$w_1$ 会收敛到 4，而 $w_2$ 永远不会更新，因为除最左状态外所有状态中的强调均为零。

现在考虑应用二步半梯度 TD 方法。
本章前面不带兴趣和强调的方法（(9.15)、(9.16) 以及第 209 页框中的算法）仍会收敛到 $\mathbf{w}_\infty=(3.5,1.5)$；而带兴趣和强调的方法会收敛到 $\mathbf{w}_\infty=(4,2)$。
后者为第一个状态和第三个状态（第一个状态从其自举）产生精确正确的价值，而永远不会进行对应于第二或第四状态的更新。
\end{minipage}}
\end{center}

\subsection*{9.12\quad 总结}

如果强化学习系统要适用于人工智能或大型工程应用，就必须能够\textbf{泛化}。
为此，只需把每次更新当作训练样例，就可以使用大量已有的\textbf{监督学习函数近似}方法。

也许最适合的监督学习方法，是那些使用\textbf{参数化函数近似}的方法；其中策略由权重向量 $\mathbf{w}$ 参数化。
虽然权重向量有许多分量，但状态空间仍然大得多，因此我们必须接受近似解。
我们定义了均方价值误差 $\overline{\mathrm{VE}}(\mathbf{w})$，作为在同策略分布 $\mu$ 下，权重向量 $\mathbf{w}$ 对价值 $v_{\pi_{\mathbf{w}}}(s)$ 的误差度量。
$\overline{\mathrm{VE}}$ 给出了在同策略情形中比较不同价值函数近似的清晰方式。

为了找到好的权重向量，最流行的方法是随机梯度下降（SGD）的各种变体。
本章集中讨论固定策略下的同策略情形，也称为策略评估或预测。
这一情形下的自然学习算法是 $n$ 步半梯度 TD，其中包含梯度蒙特卡洛和半梯度 TD(0) 作为特殊情形，分别对应 $n=\infty$ 和 $n=1$。
半梯度 TD 方法不是真正的梯度方法。
在这类自举方法（包括 DP）中，权重向量出现在更新目标中，但计算梯度时并未考虑这一点，因此它们是半梯度方法。
因此，它们不能依赖经典 SGD 结果。

尽管如此，在线性函数近似的特殊情形中，半梯度方法能够得到良好结果；线性函数近似中，价值估计是特征乘以相应权重后的求和。
线性情形在理论上理解最充分，并且在给定合适特征时实践效果良好。
选择特征是向强化学习系统加入先验领域知识的最重要方式之一。
特征可以选为多项式，但在强化学习通常考虑的在线学习情形中，多项式泛化较差。
更好的选择是傅里叶基，或者某种具有稀疏重叠感受野的粗编码。
瓦片编码是一种特别计算高效且灵活的粗编码形式。
当平滑变化的响应很重要时，径向基函数适用于一维或二维任务。
LSTD 是数据效率最高的线性 TD 预测方法，但所需计算与权重数的平方成正比；其他所有方法的复杂度都与权重数线性相关。
非线性方法包括由反向传播和 SGD 变体训练的人工神经网络；近年来，这些方法以深度强化学习的名义变得非常流行。

线性半梯度 $n$ 步 TD 在标准条件下保证对所有 $n$ 收敛到一个 $\overline{\mathrm{VE}}$，该值处于最优误差（蒙特卡洛方法渐近达到的误差）的一个界内。
这个界随着 $n$ 增大总是更紧，并在 $n\to\infty$ 时趋于零。
然而，实践中很大的 $n$ 会导致学习非常慢，因此通常偏好某种程度的自举（$n<\infty$），正如我们在第 7 章表格型 $n$ 步方法的比较，以及第 6 章表格型 TD 与蒙特卡洛方法的比较中看到的那样。

\section*{书目与历史评论}

泛化和函数近似一直是强化学习不可分割的一部分。
Bertsekas 和 Tsitsiklis（1996）、Bertsekas（2012）以及 Sugiyama 等（2013）给出了强化学习中函数近似技术现状的介绍。
强化学习中一些早期函数近似工作将在本节末尾讨论。

\noindent\textbf{9.3}
用于最小化监督学习均方误差的梯度下降方法众所周知。
Widrow 和 Hoff（1960）提出了最小均方（LMS）算法，它是典型的增量梯度下降算法。
许多教材给出了该算法及相关算法的细节（例如 Widrow and Stearns, 1985；Bishop, 1995；Duda and Hart, 1973）。

半梯度 TD(0) 最早由 Sutton（1984, 1988）作为线性 TD($\lambda$) 算法的一部分探索；我们将在第 12 章处理该算法。
用“半梯度”描述这些自举方法，是本书第二版的新术语。

强化学习中最早使用状态聚合的工作，可能是 Michie 和 Chambers 的 BOXES 系统（1968）。
强化学习中状态聚合的理论由 Singh、Jaakkola 和 Jordan（1995）以及 Tsitsiklis 和 Van Roy（1996）发展。
从动态规划最早期开始，状态聚合就已用于动态规划（例如 Bellman, 1957a）。

\noindent\textbf{9.4}
Sutton（1988）证明了在线性无关特征向量 $\{\mathbf{x}(s):s\in\mathcal{S}\}$ 情形下，线性 TD(0) 在均值意义上收敛到最小 $\overline{\mathrm{VE}}$ 解。
大约同一时期，多位研究者证明了以概率 1 收敛（Peng, 1993；Dayan and Sejnowski, 1994；Tsitsiklis, 1994；Gurvits, Lin, and Hanson, 1994）。
此外，Jaakkola、Jordan 和 Singh（1994）证明了在线更新下的收敛。
所有这些结果都假设特征向量线性无关，这意味着 $\mathbf{w}_t$ 的分量至少和状态一样多。
更重要的一般（相关）特征向量情形的收敛，最早由 Dayan（1992）证明。
Tsitsiklis 和 Van Roy（1997）显著推广并加强了 Dayan 的结果。
他们证明了本节给出的主要结果，即线性自举方法渐近误差的界。

\noindent\textbf{9.5}
我们对线性函数近似各种可能性的介绍，基于 Barto（1990）。

\noindent\textbf{9.5.2}
Konidaris、Osentoski 和 Thomas（2011）引入了一种简单形式的傅里叶基，适合具有多维连续状态空间且函数不必周期的强化学习问题。

\noindent\textbf{9.5.3}
coarse coding 这一术语来自 Hinton（1984），我们的图 9.6 基于他的图之一。
Waltz 和 Fu（1965）提供了这类函数近似在强化学习系统中的早期例子。

\noindent\textbf{9.5.4}
瓦片编码，包括哈希，由 Albus（1971, 1981）提出。
他用其“小脑模型关节控制器”（CMAC）来描述它，在文献中瓦片编码有时也以 CMAC 为名。
“tile coding”这一术语是本书第一版的新术语，不过用这些术语描述 CMAC 的想法来自 Watkins（1989）。
瓦片编码已用于许多强化学习系统，也用于其他类型的学习控制系统。
本节大量借鉴 Miller 和 Glanz（1996）的工作。
通用瓦片编码软件有多种语言实现，例如见 \texttt{http://incompleteideas.net/tiles/tiles3.html}。

\noindent\textbf{9.5.5}
自 Broomhead 和 Lowe（1988）把径向基函数与 ANN 联系起来以后，使用 RBF 的函数近似就受到广泛关注。
Powell（1987）回顾了 RBF 的早期用途，Poggio 和 Girosi（1989, 1990）则广泛发展并应用了这种方法。

\noindent\textbf{9.6}
自动调整步长参数的方法包括 RMSprop（Tieleman and Hinton, 2012）、Adam（Kingma and Ba, 2015）、随机元下降方法（例如 Delta-Bar-Delta，Jacobs, 1988）、其增量推广以及非线性推广。
专为强化学习设计的方法包括 AlphaBound、SID 和 NOSID、TIDBD，以及随机元下降在策略梯度学习中的应用。

\noindent\textbf{9.6}
McCulloch 和 Pitts（1943）把阈值逻辑单元作为抽象模型神经元提出，这是 ANN 的开端。
ANN 作为分类或回归学习方法的历史大致经历了几个阶段：单层 ANN 的 Perceptron 和 ADALINE 阶段，多层 ANN 的误差反向传播阶段，以及当前强调表示学习的深度学习阶段。
关于 ANN 的教材包括 Haykin（1994）、Bishop（1995）和 Ripley（2007）。

ANN 作为强化学习中的函数近似器，可以追溯到 Farley 和 Clark（1954）的早期工作。
此后，许多工作把类强化学习、actor--critic、反向传播、REINFORCE、TD-Gammon、AlphaGo、AlphaGo Zero 和 AlphaZero 等方法与多层 ANN 或深度卷积 ANN 联系起来。
Schmidhuber（2015）回顾了 ANN 在强化学习中的应用，包括循环 ANN 的应用。

\noindent\textbf{9.8}
LSTD 来自 Bradtke 和 Barto，并由 Boyan、Nedić 和 Bertsekas、Yu 等进一步发展。
逆矩阵的增量更新至少自 Sherman 和 Morrison（1949）以来已为人所知。
最小二乘方法向控制的扩展由 Lagoudakis 和 Parr（2003）提出。

\noindent\textbf{9.9}
我们对基于记忆函数近似的讨论，主要基于 Atkeson、Moore 和 Schaal（1997）关于局部加权学习的综述。
Atkeson（1992）讨论了基于记忆的机器人学习中局部加权回归的使用，并提供了涵盖该思想历史的大量书目。
许多后续工作把基于记忆的方法、局部加权回归、最近邻方法以及 $k$-d 树用于强化学习或相关学习控制问题。

\noindent\textbf{9.10}
核回归起源于 Aizerman、Braverman 和 Rozonoer（1964）的势函数方法。
他们把数据类比为分布在空间中、具有不同符号和大小的点电荷。
对空间中由点电荷势能求和得到的电势，对应于插值曲面。
在这个类比中，核函数就是点电荷的势能，它随到电荷的距离倒数衰减。
Connell 和 Utgoff（1987）把 actor--critic 方法应用于杆平衡任务，其中 critic 通过逆距离加权的核回归近似价值函数。
其他基于核的强化学习方法还包括 Ormoneit and Sen（2002）、Dietterich and Wang（2002）、Xu 等（2005）、Taylor and Parr（2009）、Barreto、Precup and Pineau（2011）以及 Bhat、Farias and Moallemi（2012）。

\noindent\textbf{9.11}
关于 Emphatic-TD 方法，见第 11.8 节的书目注释。

我们所知最早将函数近似方法用于学习价值函数的例子，是 Samuel 的跳棋程序（1959, 1967）。
Samuel 遵循 Shannon（1950）的建议，即价值函数不必精确到能够作为博弈中选择走法的有用指南，并且可以用特征的线性组合近似。
除了线性函数近似，Samuel 还尝试了查找表以及称为签名表的分层查找表。

大约与 Samuel 工作同时，Bellman 和 Dreyfus（1959）提出在 DP 中使用函数近似方法。
如今，函数近似方法与 DP 已有相当广泛的文献，例如多重网格方法，以及使用样条和正交多项式的方法。

Holland（1986）的分类器系统使用选择性特征匹配技术，在状态--动作对之间泛化评价信息。
每个分类器匹配一组状态，这些状态在一部分特征上具有指定取值，其余特征可以取任意值（“通配符”）。
这些子集随后被用于传统状态聚合式的函数近似。
Holland 的思想是使用遗传算法演化一组分类器，使它们共同实现有用的动作价值函数。
Holland 的思想影响了作者们早期关于强化学习的研究，但我们后来转向不同的函数近似方法。
作为函数近似器，分类器在若干方面受限。
它们本质上是状态聚合方法，因此在扩展性和高效表示平滑函数方面有相应限制。
此外，分类器的匹配规则只能实现与特征轴平行的聚合边界。
传统分类器系统最重要的限制也许是，分类器通过遗传算法学习，而遗传算法是一种进化方法。
如第 1 章所讨论，学习过程中可用的关于如何学习的信息，远比进化方法能够使用的更细致。
这一观点使我们转而把监督学习方法，尤其是梯度下降和 ANN 方法，改造后用于强化学习。
Holland 方法与我们方法之间的差异并不令人意外，因为 Holland 的思想形成于 ANN 通常被认为计算能力太弱而无用的时期；而我们的工作发生在这种传统看法开始被广泛质疑的时期。
这些不同方法的若干方面仍有许多结合机会。

Christensen 和 Korf（1986）在国际象棋中试验了用回归方法修改线性价值函数近似系数。
Chapman 和 Kaelbling（1991）以及 Tan（1991）改造决策树方法来学习价值函数。
基于解释的学习方法也被改造用于学习价值函数，从而产生紧凑表示（Yee, Saxena, Utgoff, and Barto, 1990；Dietterich and Flann, 1995）。

\end{document}
